% Source-derived chapter 07: Comparing rubber and log spaces
% Original source: pixtons-formula-abel-jacobi-picard-stack
\chapter{Comparing rubber and log spaces}
\label{pixtons-formula-abel-jacobi-picard-stack-chapter-07}
\source{pixtons-formula-abel-jacobi-picard-stack}

\begin{remark}[Source section]
\label{pixtons-formula-abel-jacobi-picard-stack-chapter-07-source}
\source{pixtons-formula-abel-jacobi-picard-stack}
\source{pixtons-formula-abel-jacobi-picard-stack}
This chapter reproduces the corresponding section of the retrieved
LaTeX source, including its definitions, statements, examples, and
proofs. It has not yet been translated into Lean.
\notready
\end{remark}

% The source section title was: Comparing rubber and log spaces
\label{sec:comparing_stacks}
\source{pixtons-formula-abel-jacobi-picard-stack}
%In this section we begin by working over a base scheme $\Lambda$, e.g. take $\bb Z$ or $\bb C$ (slightly more generality than elsewhere in the paper); we will introduce assumptions on the base later when they are needed. 

\subsection{Overview}
Our goal here is to compare the stack of stable rubber maps 
associated to a line bundle $\ca L$ on a target $X$   (introduced by Li \cite{Li2002A-degeneration-} and studied by Graber-Vakil \cite{Graber2005Relative-virtua}) to
the stack $\cat{Rub}_{g,A}$ of Marcus-Wise (see \cref{sec:log_def}) and our operational class $\DRop_{g,A}$. 
Rubber maps are reviewed in Section \ref{sec:prestable_rubber_maps} and connected to
the logarithmic space in Section \ref{clog}.
The relationship between  the construction
of Marcus-Wise and $\DRop_{g,A}$ is 
 \cref{lem:DRop_eq_MW} of Section \ref{kkkddd3}. The comparison to the class of Graber-Vakil is carried out in \cite{Marcus2017Logarithmic-com} in the case where the target $X$ is a point. We require the case where
$$X = \bb P^\kk  \ \ \ \text{and}\ \ \  \ca L = \ca O(1)\, ,$$
but only over the unobstructed locus 
$$\Mbar_{g,n}(\bb P^l, d)' \subset \Mbar_{g,n}(\bb P^l, d)\, ,$$
see \cref{sec:projective_target}. 
We will treat the case of a general nonsingular 
projective target $X$ since restricting to $\bb P^\kk$ provides no simplification (though the unobstructed locus may be rather small for general $X$). 
The final comparison result is Proposition \ref{prop:vfc_comparison_summary} in Section \ref{mwgv}.


\subsection{Refined definition of the logarithmic rubber space}

As described in Section \ref{sec:rub_log_def}, Marcus and Wise define $\cat{Rub}^{\mathsf{rel}}_{g}$ to be the moduli space of pairs $(C,P, \alpha)$ where $P$ is a tropical line on $S$ and $$\alpha\colon C \to P$$ is an $S$-morphism such that on each geometric fibre over $S$ the values taken by $\alpha$ on the irreducible components of $C$ are totally ordered in $(\bar{M}_S^{gp})_s$. However, with the above definition, certain key results of their paper (in particular concerning the comparison to spaces of rubber maps) are not correct as stated. 

To explain the problem, we restrict to the case where the base $C$ is a geometric log point. Subdividing $P$ at the images of the vertices of $C$ under the map $\alpha$ yields a \emph{divided tropical line} $Q$ (in the language of \cite{Marcus2017Logarithmic-com}). It is asserted in the discussion above \cite[\MWref{Proposition 5.5.2}\MWvone{Proposition 5.14}]{Marcus2017Logarithmic-com} that the fibre product $C \times_P Q$ is again a log curve over $S$, which, in general,
is not true. For example, take $\bar M_S$ to be the sub-monoid of $\mathbb Z^2$ generated by $(1,1)$, $(1,0)$, and $(1,-1)$, and $C$, $\alpha$ to be as illustrated in Figure \ref{fig:1}. In the fibre product,
the edge with length $(1,0)$ must be subdivided into two shorter edges, but $(1,0)$ 
is an irreducible element of $\bar M_S$. 
In fact, {\em failure of divisibility} is the only thing that can go wrong.

\begin{lemma}
\source{pixtons-formula-abel-jacobi-picard-stack}
\notready\label{lem:divisibility}
\source{pixtons-formula-abel-jacobi-picard-stack}
Let $(C, P, \alpha)$ be a point of $\cat{Rub}_g^{\mathsf{rel}}$ over a geometric log point $B$, and let $Q$ be obtained from $P$ by subdividing at the image of $\alpha$. Then the following are equivalent:
\begin{enumerate}
\item[(i)]
The fibre product $C \times_P Q$ is a log curve over $B$.
\item[(ii)] Let $e$ be an edge of $\Gamma_C$ between vertices $u$ and $v$ (satisfying $\alpha(v) \ge \alpha (u)$) with length $\ell_e\in \bar M_S$ and slope $\kappa_e = \frac{\alpha(v) - \alpha(u)}{\ell_e}$. Then, for every $y \in \on{Image}(\alpha)$ with $\alpha(u) < y < \alpha(v)$, the monoid $\bar M_{B}$ contains the element $\frac{y - \alpha(u)}{\kappa_e}$. 
\end{enumerate}
\end{lemma}
\begin{proof}
The characteristic monoid at a singular point with length $\ell_e$ is given by the monoid
\begin{equation}
\{(a,b) \in \bar M_b^2 : \ell_e \mid a-b\}\, . 
\end{equation}
Taking the fibre product over $P$ with $Q$ subdivides the characteristic monoid at the element $$\frac{y - \alpha(u)}{\kappa_e} \in \bar M_B^{\mathsf gp} \otimes_{\bb Z} \bb Q\, .$$
If $\frac{y - \alpha(u)}{\kappa_e}$ lies in $\bar M_B$, then the fiber product is easily seen to be a log curve. If not, then the subdivision is not even reduced. 
\end{proof}

\newcommand{\newRub}{\widetilde{\cat{Rub}}}
\newcommand{\newRubrel}{\newRub^{\mathsf{rel}}}
\begin{definition}
\source{pixtons-formula-abel-jacobi-picard-stack}
\notready
We define $\newRubrel$ to be the full subcategory of $\cat{Rub}^{\mathsf{rel}}$ consisting of objects $(C, P, \alpha)$ which, on each geometric fibre over $B$, satisfy the equivalent conditions of Lemma \ref{lem:divisibility}. We define $\newRub$ to be the fibre product of $\newRubrel$ over $\Picrel$ with $\Picabs$. 
\end{definition}

\begin{remark}
\source{pixtons-formula-abel-jacobi-picard-stack}
\notready
The double ramification cycle $\DRop$ can be defined as the operational class induced by the map $\cat{Rub} \to \Picabs$ following Definition \ref{defbivariantclass}. Applying the same definition to the composite map $\widetilde{\cat{Rub}} \to \Picabs$ yields the same operational class, by \Cref{prop:invarianceproperbirat}. 
\end{remark}

% As mentioned in Section \ref{sec:proof_of_equiv_defs}, we deduce from Proposition \ref{prop:invarianceproperbirat} that
% \begin{lemma}\label{lem:widetilde_same_class}

% Hmm, for comparing Rub and Div we used Prop \ref{prop:invarianceproperbirat}. But this required representability, so does not apply to a root stack. Actually, we might have a problem; the main DR construction in \ref{c44c} relies upon the fact that the map Div to Pic if representable; but this is not the case anymore with $\widetilde{\cat{Rub}}$. The map is relatively DM, so not so far from representable, maybe that helps? It would be REALLY nice not to have to make significant changes to Rosa's section... 
% \end{lemma}

\begin{figure}
\begin{tikzpicture}
\begin{scope}[every node/.style={circle,thick,draw}]
    \node (A) at (0,0) {};
    \node (B) at (0,6) {};
    \node (C) at (3,3) {};
    \node (D) at (9,0) {};
    \node (E) at (9,3) {};
    \node (F) at (9,6) {} ;
\end{scope}
    \node (G) at (4,3) {} ;
        \node (H) at (7,3) {} ;
     \node at (10,0) {$(0,0)$} ;
\node at (10,3) {$(1,1)$} ;
\node at (10,6) {$(2,0)$} ;
\begin{scope}[>={Stealth[black]},
              every node/.style={fill=white,circle},
              every edge/.style={draw=black ,very thick}]
    \path [-] (A) edge node {$(1,0)$} (B);
    \path [-] (B) edge node {$(1,-1)$} (C);
    \path [-] (C) edge node {$(1,1)$} (A);
  \path [->] (G) edge node {$\alpha$} (H);
   \path [-] (D) edge  (E);
      \path [-] (F) edge  (E);
\end{scope}
\end{tikzpicture}
\caption{A point of $\cat{Rub}$}\label{fig:1}
\source{pixtons-formula-abel-jacobi-picard-stack}
\end{figure}



\subsection{The stack of prestable rubber maps}\label{sec:prestable_rubber_maps}
\source{pixtons-formula-abel-jacobi-picard-stack}
Let $\frak M(X)$ be the stack of maps from marked prestable curves to $X$. An $S$-point of $\frak M(X)$ is a pair 
$$(C/S,\, f\colon C \to X)  $$ 
where $C/S$ is prestable with markings. To simplify notation, we will often
suppress the markings.
%Then $\MbarX$ comes with a natural map to $\frak  M(X)$. 

The space of rubber maps associated to a line bundle $\ca L$ on $X$ is summarised in \cite{Janda2018Double-ramifica}:
\emph{a map to rubber with target $X$ is a map to a rubber chain of $\bb C\bb P^1$-bundles $\bb P(\ca O_X \oplus \ca L)$ over $X$ attached along their 0  and $\infty$ divisors.}

\newcommand{\R}{R}
To facilitate our comparison, we begin by writing the definition explicitly. 
Let $\mathbf{P}$ denote the projective bundle $\bb P(\ca O_X \oplus \ca L)$.
The map collapsing the fibers,
\begin{equation}\label{55772}
\source{pixtons-formula-abel-jacobi-picard-stack}
\rho: \mathbf{P} \to X\ 
\end{equation}
admits two sections $r_0,r_\infty:X\to \mathbf{P}$
corresponding to $\ca O_X$ and $\ca L$ respectively.



\begin{definition}
\source{pixtons-formula-abel-jacobi-picard-stack}
\notready\label{def:rubber_target} An $(X,\ca L)$-rubber target $(R/S,\rho, r_0,r_\infty)$ is flat,
\source{pixtons-formula-abel-jacobi-picard-stack}
proper, and finitely presented
%\{I think `projective' should be `proper and finitely presented', if we want descent to work (to get a stack of these things). I think we also have to allow them to be algebraic spaces and not necesarily schemes for the same reason, though I don't think we need to emphasise this. } family
$$ R \to S$$
and a collapsing map $\rho: R \to X_S$
with two sections 
$$r_0,r_\infty: X_S \to R\, $$
satisfying the following properties:
\begin{enumerate}
    \item [(i)]
    Every geometric fiber  $R_s$ is isomorphic over $X_s$ to
    a finite chain  
    \begin{equation}\label{v477}
\source{pixtons-formula-abel-jacobi-picard-stack}
    \mathbf{P}\cup \mathbf{P} \cup \ldots \cup \mathbf{P}
    \end{equation}
    with the components attached successively along the respective
    $0$ and $\infty$ divisors. The collapsing maps \eqref{55772} on the
    components together define $$\rho_s: R_s \to X_s\,.$$
    The $0$ and $\infty$ sections of $\rho_s$ are determined
    by the $0$ section of first component  and the $\infty$ section of
    last components of the chain \eqref{v477}.
    \item[(ii)] \'Etale locally near every point $s\in S$, the data 
    of $(R/S,\rho, r_0,r_\infty)$ is pulled back from
    a versal deformation space described by Li \cite{Li2001-Stable}
    with one dimension for every component of the singular locus of \eqref{v477}. %\{I think we're looking at Li's [STABLE MORPHISMS TO SINGULAR SCHEMES AND RELATIVE STABLE MORPHISMS]? Then these standard models are constructed in section 4.1. In def 4.4 he constructs the stack $\frak Z^{rel}$ of all of them; is it equivalent to say that our thing comes by pullback from that?}
%\textcolor{brown}{Y:Yes, I think $\frak Z$ is a 'limit' of quotient stacks $[\mathbb{A}^n/\mathbb{G}_m^n]$ which is (noncanonically) isomorphic to rigidified $\frak M_{0,2}^{ss}$. This isomorphism is explain in Abramovich-Cadman-Fantechi-Wiss.}
\end{enumerate}
\end{definition}
\begin{definition}
\source{pixtons-formula-abel-jacobi-picard-stack}
\notready\label{def:prestable_rubber_maps}
\source{pixtons-formula-abel-jacobi-picard-stack}
The stack $\cat{Rub}^{\pre}(X, \ca L)$ of prestable rubber maps to $\ca L$ is a fibred category over $\frak M(X)$ whose fibre over a map $S \to \frak M(X)$ consists of 
three pieces of data: 
\begin{enumerate}
\item[(i)]
a prestable curve $\tilde C /S$ and a partial
stabilisation{\footnote{$C_S$ is not
necessarily a stable curve.}}
map $\tau\colon \tilde C \to C_S$ which is allowed to contract
genus 0 components with 2 special points,
%with $0$ or $1$ markings,
\item[(ii)] an $(X,\ca L)$-rubber target $(R/S,\rho,r_0,r_\infty)$,
%together
%with a partial stabilization map
%$$  R \rightarrow \mathbb{P}(\mathcal{O}_{X_S} \oplus \mathcal{L}_S)$$
%over $X_S$ 
%where the sections $r_0$ and $r_\infty$ are mapped to the sections of
%$$\mathbb{P}(\mathcal{O}_{X_S} \oplus \mathcal{L}_S)\to X_S$$ determined
%by $\mathcal{O}_{X_S}$ and  $\mathcal{L}_S$ respectively.


% \item[(iii)] isomorphisms $r_0^*\omega_{R/S}^\vee \isom L_S^\vee$ and
% $r_\infty^*\omega_{R/S} \isom \ca L_S^\vee$,
 %D: for the compatibility of these isos, maybe the predeformability condition includes something like a trivialisation of $r_0^*\omega_{R/S} \otimes r_\infty^*\omega_{R/S}$? 
 
 %D: do you really want to fix these? I thought that we would just want the differences to come from S? Or maybe just fix one of them? If we fix both we seem to rigidify too much... 

%such that the quotient $[\R/\bb G_m]$ descends to $S$; 
\item[(iii)] a map $\tilde{f}: \tilde C \to \R$ for which 
the following diagram commutes: 
%\{Check def of predeformable here - as a map over what base?}
\begin{equation}
\begin{tikzcd}
\tilde C \arrow[r,"\tilde{f}"] \arrow[d,"\tau"] & R \arrow[d, "\rho"]\\
C_S \arrow[r, "f"] & \, X_S\, .
\end{tikzcd}
\end{equation}
\end{enumerate}
The map $\tilde{f}$ in (iii) is finite over the singularities of $R/X_S$ and
predeformable{\footnote{See \cite{Li2001-Stable}.}}.
Moreover, over each geometric point $s\in S$, the image $\tilde{f}(\tilde{C}_s)$
meets
every component of $R_s$.

\end{definition}


%\item[(iv)] Normally we would say the relative conditions are imposed
%along the divisors 0 and $\infty$, but you havent included these
%conditions in the notation? \{Hmm, is that the part below this definition (just %before the end of this subsection) with the $A$ in? Or am I looking at the wrong thing? I'm not sure what's the correct way to write these in this setting, will think about. }

%\item[(v)] Equivalence ... can we just have the diagrams over a single base S?\{revised diagrams are below; but they're not really data/properties, so not sure if we want to put it in this list. }
%D: sure, I'll modify the stuff below. Just wanted to be extra careful the first time around. I have made a sketch of a proof of 47 below, but it depends on details of these definitions, in particular the meaning of predeformability, so details are not there yet. 



%\textcolor{brown}{Y: Stability condition (corresponding to R5 in MW) is missing.
%R:We are not imposing stablity yet.\\ Y: So we don't assume that for every irreducible component of $R$ there exists a component of $\tilde{C}$ which maps to this component?
%That is later forced by stability. We could also add it, but is that necessary? Then is Prop 48. true? We put this condition on $\cat{Rub}$.
%}
%\{In order to make prop 47 true I think we need to add this condition. Also the condition that no component of $\tilde C$ is mapped to a singular point of a fibre. But I'd like to check this again once the proof of 47 has all details in place.
%R: okay add it. About the not collapsing to the singular locus, I had planned
%to put that in the predef footnot, but you can put both of these two geometric
%conditions in (iii).
%}
%The map $\tilde C \to \R$ induces a line bundle $\ca O(\alpha)$ on $\tilde C$: this $\ca O(\alpha)$ can be defined by contracting away all components of $\R$ except that containing the first marking, then taking $\ca O_\R(D)$ where $D$ is that first marking.  \{Note that the contracting doesn't do anything! Also, by adjunction, $x_0^*\ca O(x_0) = x_0^*\omega^\vee$, so this is the normal bundle to the first marking $x_0$. }Our final datum is
%\begin{enumerate}[resume]
%\item we require that the line bundle $r_0^*\omega_{R/S} \otimes f^*\ca L$ on $X_S$ descends to $S$.  \{This is for comparison to $\cat{Div}^{rel}$; for comparison to $\cat{Div}$ we would instead fix a trivialisation of this line bundle. }
%an isomorphism $r_0^*\omega_{R/S}^\vee \isom f^*\ca L$. 
%\item we require that the line bundle $r_1^*\omega_{R/S}^\vee \otimes f^*\ca L$ on $X_S$ descends to $S$.  
%\item Analogously, tangent line at inf is $\ca L^\vee$ - make precise. I think this means that we should fix an iso of the tangent line at infinity to $\ca L^\vee$? But fixing it seems like too much data... I want to think about the tangent lines to one of these curves coming from MW 5.2.4... 
%\item For every section $e\colon S \to (C/S)^{ns}$ (the non-smooth locus), write $\tilde C$ for the partial normalisation along $e$, and write $h_0$, $h_1$ for the the sections lying over $e$, with $h_i$ mapping to the connected component containing $r_i$. Then we require that the line bundle
%$$h_i^*\omega_{\tilde C_T/T} \otimes f^*\ca L^{\otimes (-1)^i} $$
%(check sign) on $X_S$ descends to $S$. 

% ===older version of this condition:===\\
% For each connected component $T$ of the non-smooth locus of $C/S$, consider the base-change $C_T$ together with its natural section $t\colon T \to C_T$ landing in the non-smooth locus. We write $\tilde C_T$ for the partial normalisation along that section, and write $t_0, t_1\colon T \to \tilde C_T$ for the two sections, with $t_i$ landing in the same connected component as $r_i$. Then we require the data of isomorphisms $$t_i^*\omega_{\tilde C_T/T} \isom f_T^*\ca L^{\otimes (-1)^i}. $$
% \{Want to add in the data of these line bundles on $S$ also? Or could just require that the difference descends to $T$; not sure if it shoudl be more data? Do we require them (anti) effective on $T$ (cf. total ordering condition)?!

% Suppose we just require 
% $$t_i^*\omega_{\tilde C_T/T} \otimes f_T^*\ca L^{\otimes (-1)^i} $$
% to descend to $T$. Are they actually any more data anyway? I don't think so. And maybe these things help to pin down the log structure (or relevant part of it, for some notion of minimality...)...?

% Hmm, maybe if T -> S is not open then there is actually nothing to say, as I suspect the condition is just forced by the spreading out? What if S$S$ is $k[t]/t^2$ and the non-smooth locus is $\on{Spec} k$? I think again the condition should be automatic. 

% }

% The tangent lines at the internal nodes are $\ca L^\vee$, $\ca L$ in ordered pairs. OK, these tangent lines come in pairs, but do we fix isos? Also need to be careful in how to formalise this as the number of components varies as we move in $S$. 
%\item `the rubber $\bb G_m$ equivalence in on every component.' \{David does not understand this. Is it an extra condition?}

% ===Old version===\\
% The stack $\cat{Rub}^{\pre}(X, \ca L)$ of prestable rubber maps to $\ca L$ is a fibred category over $\frak M(X)$ whose fibre over a map $S \to \frak M(X)$ consists of 4 pieces of data: 
% \begin{enumerate}
% \item
% a prestable curve $\tilde C /S$ and a `stabilisation' map $\tau\colon \tilde C \to C_S$ (we do not require $C_S$ stable; we just allow the map to contract some genus 0 components with 0 or 1 markings); 
% \item a 2-marked semistable genus 0 curve $\R/X_S$ such that the quotient $[\R/\bb G_m]$ descends to $S$; 
% \item a predeformable map $\tilde C \to \R$ over $C_S$ which does not carry any component of any fibre into a node. 
% \end{enumerate}
% The map $\tilde C \to \R$ induces a line bundle $\ca O(\alpha)$ on $\tilde C$: this $\ca O(\alpha)$ can be defined by contracting away all components of $\R$ except that containing the first marking, then taking $\ca O_\R(D)$ where $D$ is that first marking. Our final datum is
% \begin{enumerate}[resume]
% \item
% an isomorphism $\ca O(\alpha)\isom f^*\ca L$. 
% \end{enumerate}

%Instead of fixing the isomorphism $\ca O(\alpha)\iso f^*\ca L$ we could just require its existence (up to line bundles pulled back from $S$). This is essentially what one would obtain by working with the relative Picard space instead of the Picard stack (as in \cite{Marcus2017Logarithmic-com}). 
%\{Actually let's stick to $\cat{Div}^{rel}$ for now, I think that will make more direct. Difference is not important, as absolute is $\bb G_m$ gerbe over relative. }

An isomorphism between two objects
\begin{equation*}
(\tilde C \to C_S, R, r_0, r_\infty, \tilde C \to R) \,\,\, \text{ and } \,\,\, (\tilde C' \to C_S, R', r'_0, r'_\infty, \tilde C' \to R')
\end{equation*}
over $S \to \frak M(X)$ is given by the data of isomorphisms
\begin{equation*}
\tilde C' \isom \tilde C% \,\,\, \text{ and } \,\,\, R' \isom R
\end{equation*}
over $C_S$ and
\begin{equation*}
R' \isom R
\end{equation*}
over $X_S$, compatible with the markings and such that the diagram
\begin{equation*}
\begin{tikzcd}
\tilde C' \arrow[r, "\sim"] \arrow[d] & \tilde C \arrow[d]\\
R' \arrow[r,"\sim"] & R
\end{tikzcd}
\end{equation*}
commutes. We leave the definition of the cartesian morphisms to the careful reader. 
% Suppose we are given a morphism $S' \to S$ over $\frak M(X)$, with objects 
% \begin{equation}
% (\tilde C \to C, R, r_0, r_\infty, \tilde C \to R) \,\,\, \text{ and } \,\,\, (\tilde C' \to C', R', r'_0, r'_\infty, \tilde C' \to R')\, .
% \end{equation}
% A cartesian morphism over $S' \to S$ consists of isomorphisms
% \begin{equation}
% \tilde C' \isom \tilde C \times_S S' \,\,\, \text{ and } \,\,\, R' \isom R \times_{X_S} X_{S'} 
% \end{equation}
% compatible with the markings and such that the diagrams
% \begin{equation}
% \begin{tikzcd}
% \tilde C' \arrow[r] \arrow[d] & \tilde C \times_S S' \arrow[d]\\
% C_{S'} \arrow[r] & C \times_S S'\, ,
% \end{tikzcd}
% \end{equation}
% \begin{equation}
% \begin{tikzcd}
% \tilde C' \arrow[r] \arrow[d] & \tilde C \times_S S' \arrow[d]\\
% R' \arrow[r] & R \times_{X_S} X_{S'}\, ,
% \end{tikzcd}
% \end{equation}
% and
% \begin{equation}
% \begin{tikzcd}
% R' \arrow[r] \arrow[d] & R \times_S S' \arrow[d]\\
% X_{S'} \arrow[r] & X_S
% \end{tikzcd}
% \end{equation}
% all commute. 




Suppose now that we fix a genus $g$ and a vector of integers $A$ of length $n$. We define the stack $\cat{Rub}^{\pre}_{g,A}(X, \ca L)$ with objects being tuples 
\begin{equation}\label{vtt5}
\source{pixtons-formula-abel-jacobi-picard-stack}
(\tau\colon (\tilde C, p_1,\ldots,p_n) \to C_S\, , \R/X_S\, , \tilde{f}\colon \tilde C \to \R)
%\ca O(\alpha)\isom f^*\ca L(-\sum_i a_i p_i))
\end{equation}
where $(\tilde{C},p_1,\ldots,p_n)$ is a prestable curve of genus $g$ with $n$
markings.
The data \eqref{vtt5} are as for $\cat{Rub}^{\pre}(X, \ca L)$. Moreover, 
\begin{itemize}

\item[$\bullet$]
if $a_i >0$,  $p_i\in \tilde{C}$ is mapped to the $0$-divisor with
 ramification degree $a_i$,
 \item[$\bullet$]
if $a_i <0$,  $p_i\in \tilde{C}$ is mapped to the $\infty$-divisor with
 ramification degree $-a_i$,
\item[$\bullet$] if $a_i=0$, $p_i\in \tilde{C}$ is mapped to the smooth locus of $R$
away from the $0$ and $\infty$-divisors.
\end{itemize}

%\{Do we care if the $p_i$ map to the smooth locus of $C_S$? R:NO not necessary}


\subsection{Comparison to the logarithmic space}\label{clog}
\source{pixtons-formula-abel-jacobi-picard-stack}
%As outlined in \cref{sec:rub_log_def}, Marcus and Wise %\cite{Marcus2017Logarithmic-com} introduce a logarithmic %modification $\cat{Rub}^\rel_{g,A}$ of the stack %$\cat{Div}^\rel_{g,A}$
%which is designed to facilitate the comparison to the spaces of %rubber maps. The natural map $$\cat{Rub}^\rel_{g,A} \to %\cat{Div}^\rel_{g,A}$$ 
%is a proper, birational, and log \'etale and is a log %monomorphism. 
%Hence, $\cat{Rub}_{g,A}$ and $\cat{Div}_{g,A}$ define the same %operational Chow class, see \cref{prop:invarianceproperbirat}.
%We only need to concern ourselves with comparing %$\cat{Rub}_{g,A}$ to the stack of rubber maps of %\cref{def:prestable_rubber_maps}. 


%We want to show that the space above corresponds to the universal DR $\on{Div}_0 \to \on{Pic}$, where $\on{Pic}$ is the Picard space of the universal prestable curve over $\frak M$. Marcus and Wise introduce a log modification $\on{Rub}_0 \to \on{Div}_0$, and I think this is much closer to Li's space of rubber maps. Because $\on{Rub}_0$ is a log modification of $\on{Div}_0$, they should have the same operational Chow class [check! Rosa looking in to this :-) ], so it should not matter which we use in the comparison. 


%, the open substack $\cat{Rub}_A \hra \cat{Rub}_{n}$ is defined by pulling back $\cat{Div}_A$. 
The pullback of $\ca L$ from $X$ to the universal curve over $\frak M(X)$ induces a map $\frak M(X) \to \frak{Pic}$. The key comparison result is the following. 

%Note that there is a map $\cat{Rub}^{rel} \to \Picabs$ where we rigidify the tropical line $\ca P$ by requiring that the minimum of the images of the irreducible components is given by $1$. \{I think this map is not proper; take care. Also look out for the dimension . Something is not right still. }


\begin{proposition}
\source{pixtons-formula-abel-jacobi-picard-stack}
\notready\label{lem:prestable_stack_comparison}
\source{pixtons-formula-abel-jacobi-picard-stack}
The stack $\cat{Rub}_{g,A}^{\pre}(X, \ca L)$ is naturally isomorphic to the fibre product of $\newRub_{g,A}$ over $\Picabs$ with $\frak M(X)$ along the map induced by $\ca L$,
$$\cat{Rub}_{g,A}^{\pre}(X, \ca L)\,  \stackrel{\sim}{=}\,  \newRub_{g,A} \times_{\Picabs} {\frak{M}(X)} \, .$$
%Similarly, $\frak M_A(\ca L)$  is the fibred product of $\on{Rub}_A$ over $\frak{Pic}$ with $\frak M(X)$ along the map induced by $\ca L$; or equivalently the fibred product of $\on{Rub}$ over $\frak{Pic}$ with $\frak M(X)$ along the map induced by $\ca L(-AP)$
\end{proposition}
%We could equivalently describe it as the fibred product of $\cat{Rub}_{g,\bd 0}$ over $\Picabs$ with $\frak M(X)$ along the map induced by $\ca L(-AP)$, but this is not the perspective needed for our later comparison of the virtual classes. 
% \begin{proof}
% The isomorphism is almost immediate from the discussion above \cite[\MWref{Proposition 5.5.2}\MWvone{Proposition 5.14}]{Marcus2017Logarithmic-com}. The only significant difference between $\cat{Rub}_{g,A}^{\pre}(X, \ca L)$ and the stack of Marcus and Wise is that we have a predeformability condition, and Marcus and Wise work with logarithmic curves. However, after equipping our curves with the minimal log structures, predeformability of $\tilde C \to R$ is equivalent to the map lifting to a map of log curves, by \cite{BumsigKim2008LOGARITHMICSTABLEMAPS}. %We also work with $\Picabs$ instead of $\Picrel$, but the only difference there is whether we fix an isomorphism. 
% \end{proof}
\begin{proof}
The right hand side comes with a built-in log structure, but the left side
does not. Our isomorphism will be between the underlying stacks. Our proof is based on the discussion above \cite[\MWref{Proposition 5.5.2}\MWvone{Proposition 5.14}]{Marcus2017Logarithmic-com}, and we will use the language of \emph{divided tropical lines} of \cite{Marcus2017Logarithmic-com}.

We begin by building a map from the right to the left. We are given a log curve $C/S$, a tropical line $\ca P$ on $S$, a map $\alpha\colon C\to \ca P$ whose image is totally ordered, and a map $f\colon C \to X$, such that $f^*\ca L$ lies in the isomorphism class $\ca O_C(\alpha)$. 

The $\bb G_m^{\mathsf{trop}}$-torsor $\ca P$ is rigidified by the least element among the images of the irreducible components of $C$ (here we use the total ordering condition), and hence comes with a canonical $\bb G_m$-torsor $P \to \ca P$ (if we use the rigidification to identify $\ca P = \bb G_m^{\mathsf{trop}}$ then $P = \bb G_m^{\mathsf{log}}$). The pullback $\alpha^*P$ gives a \emph{canonical} $\bb G_m$-torsor on $C$, which is isomorphic to $f^*\ca L^*$ up to pullback from $S$. In other words, the bundle $\alpha^*P \otimes f^*\ca L^\vee$ descends to a line bundle on $S$ which we denote $\ca M$. 
% The latter condition can be rephrased as saying that we have a $\bb G_m$-equivariant map $$h\colon f^*\ca L \to \ca P$$
% over $C$ (here $\bb G_m$ acts trivially on $\ca P$). 
%and an isomorphism $f^*\ca L \to \ca O_C(\alpha)$ on $C$. 

The images of the irreducible components of $C$ yield a subdivision $\ca Q$ of $\ca P$, and we define a destabilisation $\tilde C = C \times_{\ca P}\ca Q$ of $C$, which is a log curve over $S$ by Lemma \ref{lem:divisibility}. This $\ca Q$ comes with a canonical $\bb G_m$-torsor $Q$ by pulling back $P$ from $\ca P$; this $Q$ is then a 2-marked semistable genus 0 curve by \cite[\MWref{Proposition 5.2.4}\MWvone{Proposition 5.7}]{Marcus2017Logarithmic-com}. We define an $(X, \ca L)$-rubber target $R$ over $S$ by the formula
$$ R = \on{Hom}((\ca L \otimes \ca M)^*, Q)\,. $$
Here, we pull back %\ca M% and $Q$ to $X_S$
and take $\bb G_m$-equivariant homomorphisms over $X_S$. 

% We have a $\bb G_m^{log}$ torsor $P = \ca P \otimes_{\bb G_m^{trop}}\bb G_m^{log} = \bb G_m^{log}$ on $S$, and we define $Q = P \times_{\ca P} \ca Q$, a divided tropical line over $S$.  % 2-marked genus 0 semistable curve on $S$, which carries a $\bb G_m$ action by [Theorem 1.6 of https://arxiv.org/pdf/1105.4830.pdf] \{also something about  trivialisations of the normal bundles along its two markings - maybe this is part of `predeformable'?}.  
% We then define
% $$R = Q\times_S X_S \otimes_{\bb G_m} \ca (\ca L^\vee)^*\, ,$$ 
% an $(X, \ca L)$-rubber target over $S$\{Check `Li versal deformation pullback' stuff, and should there be a dual?}. We can equivalently describe $R$ as the space of $\bb G_m$-equivariant maps from $\ca L^*$ to $Q$ over $S$. 

Write $\tilde f\colon \tilde C \to X$. 
We need a predeformable map $\tilde C \to R$, equivalently an equivariant logarithmic map $\tilde f^* {\ca L}^* \to Q$ over $X_S$. It is enough to give a map $f^*\ca L \to P$ (since then we can tensor over $\ca P$ with $\ca Q$), which reduces to writing down an element of 
\begin{equation}
\begin{split}
\on{Hom}_C(f^*(\ca L \otimes \ca M), \alpha^*P) &= 
\on{Hom}_C(f^*\ca L \otimes f^*\ca L^\vee \otimes \alpha^*P, \alpha^*P)\\
& = \on{Hom}_C(\alpha^*P, \alpha^*P), 
\end{split}
\end{equation}
which contains the identity. The scheme-theoretic map is predeformable as it comes from a logarithmic map, see \cite{BumsigKim2008LOGARITHMICSTABLEMAPS}. 

%Well, we have an equivariant logarithmic map $f^*\ca L^* \to P$ (this is just a re-phrasing of our isomorphism $f^*\ca L \to \ca O_C(\alpha)$), so tensoring both sides over $\ca P$ with $\ca Q$ gives the required logarithmic map $\tilde f^* {\ca L}^* \to Q$, yielding a predeformable map $\tilde C \to R$. 


Finally we check that no component of $\tilde C$ is mapped to a non-smooth point of $R$ and that every component is hit. The target $R$ is  constructed by subdividing $f^*\ca L$ at images of components of $C$, and then $\tilde C$ is constructed by subdividing $C$ at points lying over these divisions, so both assertions are clear. %; in particular, no component of $\tilde C$ will map to a non-smooth point of $R$. 

%OK, we have a map right to left! Now we have to show it is an iso. We can work over strictly local base $S$ if we like. But we have to be cunning with log structures? 
Now we construct a map from left to right. Given a prestable rubber map to $\ca L$ over a base $S$, we first need to equip $S$ with a suitable log structure. 

The curve $R/X_S$ is a map $X_S \to \frak M_{0,2}^{ss}$, giving a (minimal) log structure on $X_S$ by pullback.
%\Ycomment{Why is it constant along $S$? Isn't it $X_S \to \frak M_{0,2}^{ss}$}\D: Yes, thanks! 
Lemma \ref{bbll} below shows that this log structure descends to $S$. The curve $R/X_S$ now carries the structure of a log curve, and similarly the quotient $[R/\bb G_m]$ descends to $S$ (again by the Lemma \ref{bbll}), determining our tropical line $\ca P$ --- which evidently satisfies the divisibility condition in Lemma \ref{lem:divisibility}.  

It remains to verify that the map $\tilde C \to R$ descends to a map $C \to \ca P$ 
and that the total ordering condition is satisfied. Write
$$\tau\colon \tilde C \to C\, .$$ 
By the proof of \cite[\MWref{Proposition 5.5.2}\MWvone{Proposition 5.14}]{Marcus2017Logarithmic-com},
we see $R \to \tau_*\tau^*R$ is an isomorphism, hence the map descends as required. The condition that no components are mapped to the nodes implies that the values of $\alpha$ on the irreducible of $C$ are a subset of the irreducible components of $R$, in particular are totally ordered. 
\end{proof}

% \begin{proof}
% The right hand side comes with a built-in log structure, but the left side
% does not. Our isomorphism will be between the underlying stacks. Our proof is based on the discussion above \cite[\MWref{Proposition 5.5.2}\MWvone{Proposition 5.14}]{Marcus2017Logarithmic-com}, and we will use the language of \emph{divided tropical lines} of \cite{Marcus2017Logarithmic-com}.

% We begin by building a map from the right to the left. We are given a log curve $C/S$, a tropical line $\ca P = \bb G_m^{\mathsf{trop}}$ on $S$, a map $\alpha\colon C\to \ca P$ whose image is totally ordered, and a map $f\colon C \to X$, such that $f^*\ca L$ lies in the isomorphism class $\ca O_C(\alpha)$. The latter condition can be rephrased as saying that we have a $\bb G_m$-equivariant map $$h\colon f^*\ca L \to \ca P$$
% over $C$ (here $\bb G_m$ acts trivially on $\ca P$). 
% %and an isomorphism $f^*\ca L \to \ca O_C(\alpha)$ on $C$. 
% %\{Nonsense. We have such a map just by projecting to $C$ then applying $\alpha$, so this cannot be a condition. I would guess that the condition is that $\alpha$ can be lifted (non-uniquely, of course) to a map $f^*\ca L^* \to \bb G_m^{log}$ over $\alpha$. Also, we work over absolute Pic, so actually the data of such a lift IS given. Not really - remember absolute Rub is just defined by pulling back relative Rub to Picabs. }

% The images of the irreducible components of $C$ yield a subdivision $\ca Q$ of $\ca P$, and we define a destabilisation $\tilde C = C \times_{\ca P}\ca Q$ of $C$. We define an $(X, \ca L)$-rubber target $R$ over $S$ by the formula
% $$ R = \on{Hom}(\ca L^*, \ca Q)\,. $$
% Here, we pull back $\ca Q$ to $X_S$
% and take $\bb G_m$-equivariant homomorphisms over $X_S$. 


% % We have a $\bb G_m^{log}$ torsor $P = \ca P \otimes_{\bb G_m^{trop}}\bb G_m^{log} = \bb G_m^{log}$ on $S$, and we define $Q = P \times_{\ca P} \ca Q$, a divided tropical line over $S$.  % 2-marked genus 0 semistable curve on $S$, which carries a $\bb G_m$ action by [Theorem 1.6 of https://arxiv.org/pdf/1105.4830.pdf] \{also something about  trivialisations of the normal bundles along its two markings - maybe this is part of `predeformable'?}.  
% % We then define
% % $$R = Q\times_S X_S \otimes_{\bb G_m} \ca (\ca L^\vee)^*\, ,$$ 
% % an $(X, \ca L)$-rubber target over $S$\{Check `Li versal deformation pullback' stuff, and should there be a dual?}. We can equivalently describe $R$ as the space of $\bb G_m$-equivariant maps from $\ca L^*$ to $Q$ over $S$. 

% Write $\tilde f\colon \tilde C \to X$. 
% We need a predeformable map $\tilde C \to R$, equivalently an equivariant logarithmic map $\tilde f^* {\ca L}^* \to \ca Q$ over $S$. But we already have an equivariant logarithmic map $h\colon f^*\ca L \to \ca P$ over $C$, and tensoring over $\ca P$ with $\ca Q$ yields the required map. The scheme-theoretic map is predeformable as it comes from a logarithmic map, see \cite{BumsigKim2008LOGARITHMICSTABLEMAPS}. 

% %Well, we have an equivariant logarithmic map $f^*\ca L^* \to P$ (this is just a re-phrasing of our isomorphism $f^*\ca L \to \ca O_C(\alpha)$), so tensoring both sides over $\ca P$ with $\ca Q$ gives the required logarithmic map $\tilde f^* {\ca L}^* \to Q$, yielding a predeformable map $\tilde C \to R$. 


% Finally we check that no component of $\tilde C$ is mapped to a non-smooth point of $R$ and that every component is hit. The target $R$ is  constructed by subdividing $f^*\ca L$ at images of components of $C$, and then $\tilde C$ is constructed by subdividing $C$ at points lying over these divisions, so both assertions are clear. %; in particular, no component of $\tilde C$ will map to a non-smooth point of $R$. 

% %OK, we have a map right to left! Now we have to show it is an iso. We can work over strictly local base $S$ if we like. But we have to be cunning with log structures? 
% Now we construct a map from left to right. Given a prestable rubber map to $\ca L$ over a base $S$, we first need to equip $S$ with a suitable log structure. 

% The curve $R/X_S$ is a map $X_S \to \frak M_{0,2}^{ss}$, giving a (minimal) log structure on $X_S$ by pullback.
% %\Ycomment{Why is it constant along $S$? Isn't it $X_S \to \frak M_{0,2}^{ss}$}\D: Yes, thanks! 
% Lemma \ref{bbll} below shows that this log structure descends to $S$. The curve $R/X_S$ now carries the structure of a log curve, and similarly the quotient $[R/\bb G_m]$ descends to $S$ (again by the Lemma \ref{bbll}) -- which determines our tropical line $\ca P$. 

% It remains to verify that the map $\tilde C \to R$ descends to a map $C \to \ca P$ 
% and that the total ordering condition is satisfied. Write
% $$\tau\colon \tilde C \to C\, .$$ 
% By the proof of \cite[\MWref{Proposition 5.5.2}\MWvone{Proposition 5.14}]{Marcus2017Logarithmic-com},
% we see $R \to \tau_*\tau^*R$ is an isomorphism, hence the map descends as required. The condition that no components are mapped to the nodes implies that the values of $\alpha$ on the irreducible of $C$ are a subset of the irreducible components of $R$, in particular are totally ordered. 
% \end{proof}
\begin{lemma}
\source{pixtons-formula-abel-jacobi-picard-stack}
\notready\label{bbll}
\source{pixtons-formula-abel-jacobi-picard-stack}
Let $(R/S, \rho, r_0, r_\infty)$ be an $(X, \ca L)$-rubber target. Then there exists a (minimal) log structure on $S$ such that $R/X_S$ can be equipped with the structure of  a log curve making $X_S$ strict over $S$. The quotient log stack $[R/\bb G_m]$ descends to a divided tropical line on $S$.  
\end{lemma}
\begin{proof}
The curve $R/X_S$ with markings $r_i$ is prestable and hence admits a minimal log structure. We must verify that the resulting log structure on $X_S$ descends to $S$. After a finite extension of $\field$ we may assume $X$ has a $\field$ point, so that $$\pi\colon X_S \to S$$
admits a section $x\colon S \to X_S$, and we can equip $S$ with the pullback log structure. It remains to construct an isomorphism $\pi^*x^* M_{X_S} \to M_{X_S}$. We start by building a map from left to right. %, equivalently a map $x^*M_{X_S} \to \pi_*M_{X_S}$. 

We first build a map on the level of characteristic monoids. The characteristic monoid at a geometric point $t \in X_S$ is given by $\bb N^\ell$, where $\ell$ is the length of the chain of projective lines of $R$ over $t$. Crucially, the irreducible elements of $\bb N^\ell$ come with a total order, given by proximity of the corresponding singularity to the $r_0$ marking. This rigidifies the characteristic monoid, so as we move along the fibre over $\pi(t)$ the characteristic monoids are \emph{canonically} identified. We
obtain canonical identifications
$$(\overline{x^*M_{X_S}})_{\pi(t)} \isom (\bar{M}_{X_S})_t\, ,$$
which give an isomorphism
$$\pi^*x^* \bar M_{X_S} \isom \bar{M}_{X_S}\, .$$


 To construct an isomorphism of log structures, we will use the perspective of \cite[Section 3.1]{borne-vistoli} that a log structure is a monoidal functor from the groupified characteristic monoid to the stack of line bundles. The rubber target is by definition pulled back from  Li's versal deformation spaces, so it suffices to construct our map in that setting. We can
 therefore assume that $S$ is regular and the locus of non-smooth curves is a reduced divisor in $X_S$.
 Since our map will be canonical, we may further shrink $S$ to be atomic\footnote{\cite[Definition 2.2.4]{Abram_wise_birational}.}. Then $\bar M_{X/S}$ is generated by its global sections, and there is a natural isomorphism of sheaves on $X_S$
 $$\phi\colon \bb N^\ell \isom \bar M_{X/S}$$
 where $\ell$ is the number of singular points in the fibre of $C$ over any point of $X_S$ lying over the closed stratum of $S$. Given $1 \le i \le \ell$, write $D_i$ for the Cartier divisor in $X_S$ where the singularity at distance $i$ from the first marking persists. Then $\phi$ sends the $i$th generator  of $\bb N^\ell$ to the section corresponding to the line bundle $\ca O_{X_S}(D_i)$. To build the required map of log structures $$\pi^*x^*M_{X_S} \isom M_{X_S}\, ,$$ we must construct an isomorphism
$$\pi^*x^*\ca O_{X_S}(D_i) \isom \ca O_{X_S}(D_i)\, .$$
Condition (i) of the \cref{def:rubber_target} implies that the underlying point set of $D_i$ is a union of fibres of $X_S/S$. Since $S$ is regular and $X_S$ is smooth over $S$, it follows that
$$D_i = \pi^*x^*D_i$$
giving the required isomorphism. %That our map of log structures is an isomorphism can easily be checked on stalks. 

% If $D_i \subseteq S$ is the locus where the  and the identification takes a singular point at distaance $i$ from the first marking to the $i$-th generator of $\bb N^r$. 

 
%  Over the closed point $s \in S$, each primitive element of the characteristic monoid lifts to a section of the characteristic monoid of the fibre, which determines a section of $\

% to which is associated a line bundle on $X_S$. We must show that these line bundles descend to $S$. Recalling that the primitive elements are totally ordered, choose some primitive element $e_i$ which is then associated to the singular point on $R_s$ at distance $i$ from the marking $r_0$. As we are in the versal setting, the associated line bundle is $\ca O_{X_S}(D_i)$ where $D_i$ is the reduced effective Cartier divisor on
% $X_S$
% %\jocomment{R: Shouldnt this be $S$ not $X$?}\{I don't think so. The curve $R$ lives over $X_S$, so once one has chosen a `crease' there is a subscheme of $X_S$ which is the locus where the crease `persists'. This will be a union of fibres over $S$. R:Ok I see,
% %I was think about the locus in S, but you are write about the inverse image in $X_S$ --- but at least it should be $X_S$ not X, right?}\{Yes. Somehow the point is to argue that the locus descends to $S$ (which is bascially obvious, but we need it scheme-theoretically, not just set-theoretically). Yes, $X_S$ indeed. Thanks! ) }
% where the $i$th singularity in the chain persists. Condition (i) of the rubber target then implies that the underlying point set of this divisor is a union of fibres of $X/S$. Since $S$ is regular 
% %\jocomment{R: Why is $S$ regular?}\{I think Li's versal deformation spaces are regular? I will check again. If not it should not be a problem - reduced will probably do. I just don't want e.g. that $D_i$ is the whole of $X_S$. Yes,they are reguar R D: Thanks again! }
% and $X_S$ is smooth over $S$, 
% $D_i$ is the pullback of a reduced effective Cartier divisor $D_i'$ on $S$. Therefore 
% $$\ca O_{X_S}(D_i) = \pi^*x^*\ca O_{X_S}(D_i)$$ as required. 


% That our map is an isomorphism can easily be checked on stalks. 

%Maybe DF pproach to log structures will make it easier? We can assume $S$ strictly local. The rank of the characteristic is constant on each fibre, in particular on the closed fibre say it is $n$, then I want to put the monoid $\bb N^n$ on $S$, and map generators to line bundles on $S$. Actualyl I will map them to line bundles on $X_S$ and observe that these line bundles (and maps between them) descend to $S$. Well, each generator corresponds to an intersection of two components of the chain, and this in turn gives a divisor on the versal deformmation space where the singularity persists. The associated line bundle is exactly the thing I want, so I just need to show that this divisor is pulled back from a suitable versal deformation of $S$ (indeed, enough to treat universal case to avoid this versal nhd annoyance? But want also $S$ strictly local; can do that after switching to universal, to get best of both worlds! ). Now in this versal setting the spaces are reduced and the divisors are honest reduced effective Cartier divisors, hence descend o reduced effective Cartier divisors on $S$ as required. 

The quotient log stack $[R/\bb G_m]$ is a divided tropical line on $X_S$ with divisions coming from the divisors $D_i$. We can identify the underlying tropical line with $\bb G_m^{trop}$ by specifying that the smallest element in the sequence of divisions is mapped to $0$. We have already established that these divisions $D_i$ descend to $S$, hence so does the divided tropical line. 
\end{proof}
%\{This seems to give an alternative description of the space of $(X, \ca L)$ rubber targets over $S$. Namely, first take a 2-marked semistable genus 0 curve $T$ over $S$, with its built-in $\bb G_m$-action. Pull this back to $T_X/X$, then take $(T_X \times_X \ca L^*)/\bb G_m$, i.e. twist the pullback by $\ca L$. Is this worth remarking? Not sure, but feels a bit weird to leave it implicit. }


% ===old===\\

%Let $f:X \to S$ be a flat map to a reduced target, and let $D$ a reduced effective Cartier divisor on $X$ such that for every geometric point $s$ of $S$, the fibre $D_s$ is either empty or equal to $X_s$. Then there exists a divisor $E$ on $S$ which pulls back to $D$. Easy. 


After restriction to the locus 
where the infinitesimal automorphisms are trivial, we obtain
 a stable version of \cref{lem:prestable_stack_comparison}. 
Let $\MbarX$ denote the stack of stable maps from $n$-pointed curves to $X$ 
representing  the class $\beta$.
The line bundle $\ca L$ determines 
a map $$\MbarX \to \Picabs\, $$  
and we can pullback $\newRub_{g,A}$ as before.
Let $$\cat{Rub}_{g,A}(X, \ca L) \subset \cat{Rub}_{g,A}^{\pre}(X, \ca L)$$ 
be the locus where the infinitesimal automorphisms are trivial.


\begin{lemma}
\source{pixtons-formula-abel-jacobi-picard-stack}
\notready
The stack $\cat{Rub}_{g,A}(X, \ca L)$ is the fibre product of $\newRub_{g,A}$ over $\Picabs$ with $\MbarX$ along the map given by $\ca L$. %, or equivalently the fibred product of $\on{Rub}_A$ over $\frak{Pic}$ with $\Mbar_n(X)$ along the map given by $\ca L$,
$$\cat{Rub}_{g,A}(X, \ca L)\,  \stackrel{\sim}{=}\,  \newRub_{g,A} \times _\Picabs \MbarX\, .$$
\end{lemma}

Next, we will compare the virtual fundamental classes on these
spaces. We will carry out the comparison on a smaller open locus.  
%see \cref{sec:projective_target}. 
We define 
\begin{enumerate}
\item[(i)]
$\MbarX'$ is be the open locus of maps
$(C, f\colon C \to X)$ in $\MbarX$ where $H^1(C, f^*\ca L) = 0$.
\item[(ii)]
$\cat{Rub}_{g,A}(X, \ca L)' = \cat{Rub}_{g,A}(X, \ca L) \times_{\MbarX} \MbarX'$. 
\end{enumerate}
In  \cref{sec:projective_target}, we considered the case $X = \bb P^\kk$
and showed that this unobstructed locus is large enough
to control the cycles relevant to Theorem \ref{Thm:main}. For general $X$,
the unobstructed  locus might be very small (and possibly empty). 



\subsection{Comparing the virtual classes}\label{sec:comparing_virtual_classes}
\source{pixtons-formula-abel-jacobi-picard-stack}
\subsubsection{Overview}
We begin by briefly discussing of several spaces which will be relevant in setting up the obstruction theories. Let $$\frak M_{g, n}^{ss}\subset \frak{M}_{g,n}$$
be the semistable locus (where every rational curve has at least two distinguished points).
We write $\cl T$ for the algebraic stack with log structure which parametrises tropical lines with at least one division. %The underlying algebraic stack is the locus of curves in $\frak M_{0,3}^{ss}$ with the first two markings landing on the same component. %, quotiented by $\bb G_m$ acting on that first component so as to fix the first marking and the node (or first and third markings, in the smooth case). 
There are natural maps $$\frak M_{0,2}^{ss} \to \cl T\ \ \ \text{and}\ \ \  \frak M_{0,2}^{ss} \to B\bb G_m\, ,$$
the former defined by dividing $\bb G_m^{\mathsf{trop}}$ at $1$ and at the smoothing parameters of the nodes, and the latter defined by the normal bundle at the first marking. The induced map 
\begin{equation}
\frak M_{0,2}^{ss} \to \cl T\times B\bb G_m
\end{equation}
is an isomorphism by \cite[Proposition 3.3.3]{Abramovich2013Expanded-degene}. %\cite[Theorem 5.2.5]{Marcus2017Logarithmic-com}. 

As $\cat{Rub}^{\rel}$ is the moduli stack of tuples $(C, \alpha\colon C \to \ca P)$ where $\ca P$ is a tropical line and the images of the irreducible components of $C$ are totally ordered, there is a natural map 
\begin{equation}\label{ff923}
\source{pixtons-formula-abel-jacobi-picard-stack}
\cat{Rub}^{\rel} \to \cl T
\end{equation}
sending $(C, \alpha\colon C \to \ca P)$ to the tropical line $\ca P$ with the division given by the images of the irreducible components of $C$. 
%By \eqref{ff923} and the
%definition \eqref{rubbdeff}
%of $\cat{Rub}_{g,A}$, we have
%a map
%\begin{equation}\label{vrt5}
% \cat{Rub}_{g,A}(X, \ca L)'\to \frak M_{0,2}^{ss}
% \end{equation}


We will construct a map 
 \begin{equation}\label{vrt5}
\source{pixtons-formula-abel-jacobi-picard-stack}
 \cat{Rub}_{g,A}(X, \ca L)'\to \frak M_{0,2}^{ss}
 \end{equation}
 lifting the  morphism \eqref{ff923} by the following argument. 
A point of $\cat{Rub}_{g,A}(X, \ca L)'$ is a tuple $(C, \alpha\colon C \to \ca P, f\colon C \to X)$ where $f^*\ca L$ lies in the class{\footnote{Here,
$[\ca O_C(\alpha)]$
 is an equivalence class under isomorphisms and tensoring with pullbacks from $S$.
}}
$[\ca O_C(\alpha)]$.  However, as $\ca P$ is divided, there is a unique isomorphism $\ca P \isom \bb G_m^{trop}$ where the smallest division maps to $0$. The universal $\bb G_m$ torsor $\bb G_m^{log} \to \bb G_m^{trop}$ pulls back to a well-defined $\bb G_m$-torsor $\ca O_C^*(\alpha)$ on $C$, and the difference $f^*\ca L^* \otimes_{\ca O^*_C} \ca O^*_C(-\alpha)$ descends to a $\bb G_m$-torsor on $S$ by the construction of $\cat{Rub}_{g,A}(X, \ca L)'$ as a fibre product. The $\bb G_m$-torsor on $S$ induces a map $\cat{Rub}_{g,A}(X, \ca L)' \to B \bb G_m$. Combined with the map $\cat{Rub}_{g,A}(X, \ca L)' \to \cl T$ via \eqref{ff923},
we obtain the
map \eqref{vrt5}.
%$\cat{Rub}_{g,A}(X, \ca L)' \to \frak M_{0,2}^{ss}$. 

% There is then a natural map $\cat{Rub}_{g,A} \to \frak M_{g,n} \times \frak M_{0,2}^{ss}$, and composition yields a natural map 
% \begin{equation}
% \cat{Rub}_{g,A}(X, \ca L)' \to \frak M_{g,n} \times \frak M_{0,2}^{ss}\,. 
% \end{equation}

%There is also a natural map $\cat{Rub}_{g,A}(X, \ca L)' \to B \bb G_m$; we send a tuple $(\tau\colon \tilde C \to C_S, \R/S, \phi\colon \tilde C \to \R, p_1, \dots, p_n)$ to the line bundle on $S$ which pulls back to $f^*\ca L \otimes \ca O(\alpha)^\vee$. \{This will change if we fix the iso earlier... May be good or bad...}

The space $\cat{Rub}_{g,A}(X, \ca L)'$ carries three virtual fundamental classes by
the following three constructions:
\begin{enumerate}
\item[(i)] The class $\DRop_{g,A}(\phi_\ca L)([\MbarX'])$ obtained by applying %$\cat{Div}_{g,A}$ 
$\DRop_{g,A}$
to the (virtual) fundamental class of $\MbarX'$ via the map $\phi_\ca L$. 
\item[(ii)] The class obtained from a two-step obstruction theory described by Marcus and Wise \cite{Marcus2017Logarithmic-com} for the map $\cat{Rub}_{g,A}(X, \ca L)' \to \frak M_{g,n} \times \cl T$. 
\item[(iii)]  A class coming from a two-step obstruction theory studied by
Graber and Vakil \cite{Graber2005Relative-virtua} for the map 
%$$\cat{Rub}_{g,A}(X, \ca L)' \to %\MbarX' \times  \cl T \, .$$ 
$$\cat{Rub}_{g,A}(X, \ca L)' \to \MbarX' \times \frak M_{0,2}^{ss}\, .$$ 
\end{enumerate}
We will prove (i)-(iii) are all equal.
%\{note to self: we use \cite[3.15]{Manolache2012Pullbacks} for the claim that composing with unobstructed maps allows us to lift prefect rel obs theories. }

% \{Dimension sanity check (to delete when all makes sense): \\
% $\frak M_{0,2}^{ss}$ is open in $\frak M_{0,2}$, so has dimension -1\\
% $B\bb G_m$ has dimension -1\\
% $\cl T$ is open in $\frak M_{0,3}$ so has dimension 0. }

%There is an algebraic stack $\cl T/\Lambda$ such that $\frak M_{0,2}^{ss}  = \cl T \times_\Lambda B\bb G_m$. \todo{say more! I suspect that $\cl T$ parametrises curves in $\frak M_{0,3}^{ss}$ with the first two markings landing on the same component. }

%pullback of $\ca O(\alpha)$ along the first section of $\R$. \todo{choices here!}

% \subsubsection{The rubber obstruction theory}
% %\todo{Younghan, should I replace the $\frak M_{0,2}^{ss}$ in this subsubsection with a $\cl T$? Y. I think $\frak M_{0,2}^{ss}$ is correct.}
% Graber and Vakil \cite{Graber2005Relative-virtua} describe a perfect relative obstruction theory for the map
% \begin{equation}
% \cat{Rub}_{g,A}(X, \ca L)' \to \frak M_{g,n} \times \frak M^{ss}_{0,2}\,. 
% \end{equation}
% We have a commutative diagram 
% \begin{equation}
%  \begin{tikzcd}
%   \cat{Rub}_{g,A}(X, \ca L)' \arrow[dr] \arrow[r] & \MbarX' \times \frak M_{0,2}^{ss} \arrow[d]\\
% & \frak M_{g,n} \times \frak M_{0,2}^{ss}\, , 
% \end{tikzcd}
% \end{equation}
% where the vertical arrow is unobstructed; hence to give obstruction theories for the diagonal and horizontal maps is equivalent by \cite[3.15]{Manolache2012Pullbacks}. 

\subsubsection{$\DRop$ and the obstruction theory of Marcus-Wise}
\label{kkkddd3}
\source{pixtons-formula-abel-jacobi-picard-stack}
% The virtual fundamental class of $\MbarX'$ equals the fundamental class.  Applying the operational class of $\cat{Rub}_{g,A}$ in $\CHop^g(\Picabs_{g,n})$ to the fundamental class of $\MbarX'$ induces a class on $\cat{Rub}_{g,A}(X, \ca L)'$. 
% To compare with the class coming from the rubber theory we will make this more explicit. 

The obstruction theory of Marcus-Wise is a \emph{two-step} obstruction theory, a notion which we now recall. Unless otherwise stated, by \emph{perfect obstruction theory} we mean an obstruction theory which is perfect in amplitude $[-1,0]$. 

\begin{definition}
\source{pixtons-formula-abel-jacobi-picard-stack}
\notreadyA \emph{two-step} obstruction theory for a map $f\colon X \to S$ consists of a factorisation $$X \to Y \to S$$ together with 
prefect relative obstruction theories for $X/Y$ and for $Y/S$. 
\end{definition}
A two-step obstruction theory induces a virtual pullback by composition.{\footnote{A
two-step obstruction theory
also induces a perfect obstruction theory for $X/S$ in amplitude $[-2,0]$, but we
will not use the latter construction.}}
If $S$ has a fundamental class $[S]$, the virtual pullback of $[S]$
is the {\em virtual fundamental class} of $X$ associated to the
two-step obstruction theory.


We first recall the two-step obstruction theory of \cite{Marcus2017Logarithmic-com} in the case when $X$ is a point. We have a diagram
\begin{equation}
 \begin{tikzcd}
  \cat{Rub}_{g,A}(\mathsf{pt}, \ca O)' \arrow[dr] \arrow[r] & \widetilde{\cat{Rub}}_{g,A} \arrow[d]\\
& \frak M_{g,n} \times \cl T. 
\end{tikzcd}
\end{equation}
A perfect relative obstruction theory for the horizontal map is given in \cite[\MWref{Section 5.6.3}\MWvone{Section 5.6.2}]{Marcus2017Logarithmic-com}, and for the vertical map in \cite[Proposition \MWref{5.6.5.3}\MWvone{5.20}]{Marcus2017Logarithmic-com}; while the reader might expect that these arguments apply to $\cat{Rub}_{g,A}$ rather than the root stack $\widetilde{\cat{Rub}}_{g, A}$, Marcus and Wise in fact assume in both constructions the divisibility conditions of Lemma \ref{lem:divisibility}, hence their constructions in fact apply to $\widetilde{\cat{Rub}}_{g, A}$ (and not to $\cat{Rub}_{g, A}$). This two-step obstruction theory
%these yields a perfect relative obstruction theory for the diagonal map. A priori, the amplitude is $[-2, 0]$, but in fact the diagonal theory 
coincides with the rubber theory of Graber-Vakil, as shown in \cite[Section \MWref{5.6.6}\MWvone{5.6.5}]{Marcus2017Logarithmic-com}.
Moreover, the virtual fundamental class obtained
equals the operational class of $\cat{Div}_{g,A}$, see \cite[\MWref{Theorem 5.6.1}\MWvone{Theorem 5.15}]{Marcus2017Logarithmic-com}; again, these results all assume the divisibility condition of Lemma \ref{lem:divisibility}, and hence apply to $\widetilde{\cat{Rub}}_{g,A}$ in place of $\cat{Rub}_{g, A}$. 

%\{GV have target $\frak M_{g,n} \times \frak M_{0,2}^{ss}$. The difference is just a $\bb G_m$-gerbe so unobstructed, but be careful. }

Returning to the case of arbitrary $(X,\ca L)$,
we can construct a similar commutative diagram 
\begin{equation}\label{eq:obs_factor_MW}
\source{pixtons-formula-abel-jacobi-picard-stack}
 \begin{tikzcd}
  \cat{Rub}_{g,A}(X, \ca L)' \arrow[dr] \arrow[r] & \widetilde{\cat{Rub}}_{g,A} \arrow[d]\\
& \frak M_{g,n} \times \cl T\, .
\end{tikzcd}
\end{equation}
The vertical map is unchanged and so again has a perfect relative obstruction theory by \cite[Proposition \MWref{5.6.5.3}\MWvone{5.20}]{Marcus2017Logarithmic-com}; in fact the morphism is a local complete intersection, and the obstruction theory of \cite[Proposition \MWref{5.6.5.3}\MWvone{5.20}]{Marcus2017Logarithmic-com} is just the relative tangent complex. 

We need to supply a perfect obstruction theory for the horizontal map $$\widetilde{\cat{Rub}}_{g,A}(X, \ca L)'  \to  \cat{Rub}_{g,A}\, ,$$
which we can factor as
\begin{equation}
\cat{Rub}_{g,A}(X, \ca L)'  \to \MbarX' \times_{\frak M_{g,n}}\widetilde{\cat{Rub}}_{g,A}  \to \widetilde{\cat{Rub}}_{g,A}\, . 
\end{equation}
The second map is a base change of the unobstructed map $\MbarX' \to \frak M_{g,n}$, hence is unobstructed. For the first map, consider the pullback square
\begin{equation}\label{eq:long_vertical_pullback}
\source{pixtons-formula-abel-jacobi-picard-stack}
 \begin{tikzcd}
\cat{Rub}_{g,A}(X, \ca L)'  \arrow[d] \arrow[r]   & \MbarX' \arrow[d]\\
\MbarX' \times_{\frak M_{g,n}} \widetilde{\cat{Rub}}_{g,A} \arrow[r]  & \MbarX' \times_{\frak M_{g,n}}\Picabs_{g,n}\, .\\
%\MbarX' \times_{\frak M_{g,n}} \cat{Rub}_{g,A} \arrow[r] \arrow[d]  & \MbarX' \times_{\frak M_{g,n}}\Picabs_{g,n}\arrow[d]\\
%\cat{Rub}_{g,A} \arrow[r]&  \Picabs_{g,n}
\end{tikzcd}
\end{equation}
The right vertical arrow is a section of a base change of the smooth morphism $\Picabs_{g,n} \to \frak M_{g,n}$, and as such is lci and has a perfect relative obstruction theory given by the relative tangent complex $R^1\pi_*\ca O_C$. Pullback yields a corresponding perfect obstruction theory for the left vertical arrow. This gives a two-step obstruction theory for the composite map $\cat{Rub}_{g,A}(X, \ca L)'  \to  \widetilde{\cat{Rub}}_{g,A}$, from which we obtain a virtual fundamental class following \cite{Manolache2012Pullbacks}. 

The discussion here
is a very slight generalisation of the obstruction theory constructed in \cite[\MWref{Proposition 5.6.3.1}\MWvone{Proposition 5.16}]{Marcus2017Logarithmic-com}. 


\begin{definition}
\source{pixtons-formula-abel-jacobi-picard-stack}
\notready
The two-step obstruction theory for the diagonal map of \eqref{eq:obs_factor_MW},
$$
\cat{Rub}_{g,A}(X, \ca L)' \to \frak M_{g,n} \times \cl T\, $$ is the \emph{Marcus-Wise} obstruction theory.
\end{definition}


% We consider the lifting problem 
% \begin{equation}
%  \begin{tikzcd}
% S \arrow[d] \arrow[r] &  \cat{Rub}_{g,A}(X, \ca L)' \arrow[d]\arrow[r] & \MbarX' \arrow[d]\\
% S' \arrow[r] \arrow[ur, dashed] \arrow[urr, dashed]& \cat{Rub}_{g,A} \arrow[r] & \Picabs_{g,n}
% \end{tikzcd}
% \end{equation}
% where $S \to S'$ is a square-zero extension. Since the right square is cartesian, it is equivalent to lift either arrow, and we focus on the outer rectangle. We can factor further, decomposing $\MbarX' \to \Picabs_{g,n}$ as the composition
% \begin{equation}
% \MbarX' \to \MbarX' \times_{\frak M_{g,n}} \Picabs_{g,n} \to \Picabs_{g,n}\,, 
% \end{equation}
% yielding a lifting problem 
% \begin{equation}\label{eq:long_vertical_lift}
%  \begin{tikzcd}
% S \arrow[dd] \arrow[rrr]&  & & \MbarX' \arrow[d]\\
% & && \MbarX' \times_{\frak M_{g,n}}\Picabs_{g,n}\arrow[d]\\
% S' \arrow[rrr] \arrow[urrr, dashed] \arrow[uurrr, dashed]& & & \Picabs_{g,n}\,.
% \end{tikzcd}
% \end{equation}
% The vertical map $\MbarX' \times_{\frak M_{g,n}}\Picabs_{g,n} \to \Picabs_{g,n}$ is unobstructed, since it is a base change of the unobstructed map $\MbarX' \to \frak M_{g,n}$ by the smooth stack $\Picabs_{g,n} /\frak M_{g,n}$, so it suffices to understand the lifting problem for the vertical map $$\MbarX' \to \MbarX' \times_{\frak M_{g,n}}\Picabs_{g,n}\, .$$

% We can make the formulation more explicit. The map $$S' \to \MbarX'\times_{\frak M_{g,n}} \Picabs_{g,n}$$ gives the data of a curve $C'/S'$, a map $f'\colon C' \to X$, and a line bundle $\ca F$ on $C'$, and the lift of $S$ gives an isomorphism 
% $$\psi\colon \ca F|_C \isom f^* \ca L = ((f')^*\ca L)|_C\, .$$ 
% To construct the dashed arrow, we need to lift the isomorphism $\psi$ to $S'$, giving an isomorphism $$\psi'\colon \ca F \iso (f')^*\ca L\, .$$
% Such lifts $\psi'$ form a torsor under
% \begin{equation}
% \on{Hom}_{C'}((f')^*\ca L, J\ca F) = \on{Hom}_C(f^*\ca L, J \otimes f^*\ca L) = J\,, 
% \end{equation}
% where $J$ is the square-zero ideal sheaf on $S'$ cutting out $S$. Hence the obstruction lies in $$H^1(C, J) = J \otimes H^1(C, \ca O_C)\, .$$
% In other words, there is a perfect relative obstruction theory for $\MbarX'$ over $\MbarX' \times_{\frak M_{g,n}} \Picabs_{g,n}$ given by $R^1\pi_*\ca O_C$. 



%One can show by an analogous argument to the proof of \cite[\MWvone{Theorem 5.15}]{Marcus2017Logarithmic-com} that the resulting virtual fundamental class coincides with that obtained by applying $\cat{Div}_{g,A}$ to the (virtual) fundamental class of $\MbarX'$. \{Should we give more explanation of this last sentence? The proof of the corresponding statement in [MW, section 5.6.2] is a pretty formal application of the log cotangent complex and Olsson's stacks of log structures, and carries over to our setting without significant change. }






\begin{lemma}
\source{pixtons-formula-abel-jacobi-picard-stack}
\notready\label{lem:DRop_eq_MW}
\source{pixtons-formula-abel-jacobi-picard-stack}
The push forward along
$$\psi\colon \cat{Rub}_{g,A}(X, \ca L)' \to \MbarX'$$
of the virtual fundamental class of the Marcus-Wise theory
on $\cat{Rub}_{g,A}(X, \ca L)'$ 
equals the class $\DRop_{g,A}(\phi_\ca L)([\MbarX'])$
obtained via the map
$$\phi_\ca L\colon \MbarX' \to \Picabs_{g,n}\, .$$
%coincides with the application of
%$\cat{Div}_{g,A}$ to the fundamental class of $\MbarX'$ via the map $\phi_\ca L$. 
\end{lemma}
\begin{proof}
%From the map $$\phi_\ca L\colon \MbarX' \to \Picabs_{g,n}\, ,$$ 
From $\phi_{\ca L}$, we obtain maps
$$\phi'_\ca L\colon  \MbarX' \to \MbarX'\times_{\frak M_{g,n}}\Picabs_{g,n}\, ,$$ $$\phi''_\ca L\colon  \MbarX' \to \MbarX'\times \Picabs_{g,n}\, . $$
Both are lci morphisms because $\Picabs_{g,n}/\frak M_{g,n}$ is smooth. By 
Definition \ref{defbivariantclass} and Section \ref{sec:proof_of_equiv_defs} we have
\begin{align*}
    \DRop_{g,A}(\phi_\ca L)([\MbarX']) &=  \psi_*(\phi''_\ca L)^![\MbarX' \times \widetilde{\cat{Rub}}_{g,A}]\\
    &= \psi_*(\phi'_\ca L)^![\MbarX' \times_{\frak M_{g,n}} \widetilde{\cat{Rub}}_{g,A}]\, .
\end{align*}
The virtual fundamental class of the Marcus-Wise theory is the virtual pullback of the fundamental class of $\frak M_{g,n} \times \cl T$ along the composition
\begin{equation}
\cat{Rub}_{g,A}(X, \ca L)'  \stackrel{1}{\to} \MbarX' \times_{\frak M_{g,n}} \widetilde{\cat{Rub}}_{g,A}   \stackrel{2}{\to} \widetilde{\cat{Rub}}_{g,A}   \stackrel{3}{\to} \frak M_{g,n} \times \cl T\, . 
\end{equation}
The map (3) is lci and the obstruction theory is the relative tangent complex, so the pullback of the fundamental class is the fundamental class of $\widetilde{\cat{Rub}}_{g,A}$. The map (2) is unobstructed, so the (virtual) pullback of the fundamental class is again the fundamental class. The obstruction theory of the map (1) is defined by pulling back the relative tangent complex of the lci morphism $$\MbarX' \to \MbarX' \times_{\frak M_{g,n}}\Picabs_{g, n}$$ via the pullback square
\begin{equation}
\begin{tikzcd}
\cat{Rub}_{g,A}(X, \ca L)' \arrow[r] \arrow[d] & \MbarX' \times_{\frak M_{g,n}} \widetilde{\cat{Rub}}_{g,A}  \arrow[d] \\
\MbarX' \arrow[r, "\phi'_\ca L"] & \MbarX' \times_{\frak M_{g,n}} \Picabs_{g,n}\, , 
\end{tikzcd}
\end{equation}
so the virtual pullback of the fundamental class of $\MbarX' \times_{\frak M_{g,n}} \widetilde{\cat{Rub}}_{g,A}$ is  equal to the Gysin pullback $(\phi'_\ca L)^![\MbarX' \times_{\frak M_{g,n}} \widetilde{\cat{Rub}}_{g,A}]$. 
% \jocomment{I dont understand the notation. For (3), arent
% you going to use the MW relative obstruction theory [55,Prop 5.20]
% Could you write a few more sentences here about how this
% gives the Marcus-Wise fundamental class of Def 51?}
% \{Yes, for (3) we use [MW prop 5.20]. But in the new version they remark (Remark 5.6.5.4.) that the moprhism is lci, and their obstruction theory is just the relative tangent complex. Hence the virtual pullback is just the gysin pullback. For (2) map is unobstructed so smooth, and we can just do easy pullback in CHow, or virtual pullback with trivial obs theory, and get same answer. For (1) argument is basicalyl same as (3); the map is base-change of lci, and the obs theory is pullback of tangent complex of that lci, so again virtual pulllback is gysin pullback. But I agree should write this,and more clearly! } 
% We must verify that the pullback of the fundamental class of $\frak M_{g,n} \times \cl T$ 
% to $\cat{Rub}_{g,A}(X, \ca L)'$
% pushes forward to $\DRop_{g,A}(\phi_\ca L)([\MbarX'])$. 
% First, the map (3) is lci by \cite[Remark 5.6.5.4]{Marcus2017Logarithmic-com}, hence the pullback of the fundamental class is the fundamental class of $\cat{Rub}_{g,A}$.
% \jocomment{The actual fundamental class of Rub?}\{Yes; Rub itself is nice and expected dimension etc, the problems only come when we start base-changing it. }
% The map (2) is unobstructed as it is the base change of the unobstructed map $$\MbarX' \to \frak M_{g,n}\, ,$$
% hence again pulling back the fundamental class yields the fundamental class. 
% For the map (1), we have a pullback square
% \begin{equation}
% \begin{tikzcd}
% \cat{Rub}_{g,A}(X, \ca L)' \arrow[r] \arrow[d] & \MbarX' \times_{\frak M_{g,n}} \cat{Rub}_{g,A}  \arrow[d] \\
% \MbarX' \arrow[r, "\phi'_\ca L"] & \MbarX' \times_{\frak M_{g,n}} \Picabs_{g,n}\, , 
% \end{tikzcd}
% \end{equation}
% so the Gysin pullback of the fundamental class of $\MbarX' \times_{\frak M_{g,n}} \cat{Rub}_{g,A}$ is by definition equal to $\DRop_{g,A}(\phi_\ca L)([\MbarX'])$. 
\end{proof}

% \subsubsection{Obstruction theory induced by Marcus-Wise}
% %{Obstructions over \texorpdfstring{$\MbarX' \times \frak M_{0,2}^{ss}$}{Mbar(X) x Mfrak\_\{0,2\}\^{}\{ss\}}}

% As recalled above, Marcus-Wise define a two-step obstruction theory for the map 
% \begin{equation*}\label{vrrt}
% \cat{Rub}_{g,A}(X, \ca L)' \to \frak M_{g,n} \times \cl T\,.
% \end{equation*}
% However, 
% Graber-Vakil consider an obstruction theory
% for the map 
% \begin{equation} \label{vrrt}
% \cat{Rub}_{g,A}(X, \ca L)' \to \MbarX' \times \frak M_{0,2}^{ss}\,.
% \end{equation}


% We show here how the Marcus-Wise obstruction theory induces an obstruction 
% theory for the map \eqref{vrrt}.
% We have a commutative diagram 
% \begin{equation}
%  \begin{tikzcd}\label{kk882}
%   \cat{Rub}_{g,A}(X, \ca L)' \arrow[dr] \arrow[r] & \MbarX' \times \frak M_{0,2}^{ss} \arrow[d]\\
% & \frak M_{g,n} \times \frak M_{0,2}^{ss}\, , 
% \end{tikzcd}
% \end{equation}
% where the vertical arrow is unobstructed. Giving an obstruction theory for the diagonal map is therefore equivalent to giving an obstruction theory for the horizontal map by \cite[3.15]{Manolache2012Pullbacks}. 
% We have $$\frak M_{0,2}^{ss} = \cl T \times B\bb G_m\, ,$$
% and the map $B \bb G_m \to \on{Spec} \field$ is unobstructed. Giving 
% an obstruction theory for the diagonal map is again equivalent to
% giving an obstruction theory for the horizontal map of 
% \begin{equation}\label{hhww33}
%  \begin{tikzcd}
%   \cat{Rub}_{g,A}(X, \ca L)' \arrow[dr] \arrow[r] & \frak M_{g,n} \times \cl T \times B\bb G_m \arrow[d]\\
% &\frak M_{g,n} \times \cl T\, . 
% \end{tikzcd}
% \end{equation}



% If we start with the Marcus-Wise obstruction theory for the diagonal map
% of \eqref{hhww33}, we obtain an obstruction theory for the horizontal map of
% \eqref{hhww33}. The latter is the diagonal map of \eqref{kk882}, so we obtain
% an obstruction theory for the horizontal map of \eqref{kk882}.
% The virtual fundamental classes of
% $\cat{Rub}_{g,A}(X, \ca L)'$
% associated to the Marcus-Wise theory and the induced obstruction theory
% of the horizontal map of \eqref{kk882} are equal by \cite{Manolache2012Pullbacks}. 

%\subsubsection{Obstruction theory induced by Marcus-Wise}



% We show here how the Marcus-Wise obstruction theory induces an obstruction 
% theory for the map \eqref{vrrt}.
% We have a commutative diagram 
% \begin{equation}
%  \begin{tikzcd}\label{kk882}
%   \cat{Rub}_{g,A}(X, \ca L)' \arrow[dr] \arrow[r] & \MbarX' \times \frak M_{0,2}^{ss} \arrow[d]\\
% & \frak M_{g,n} \times \frak M_{0,2}^{ss}\, , 
% \end{tikzcd}
% \end{equation}
% where the vertical arrow is unobstructed. Giving an obstruction theory for the diagonal map is therefore equivalent to giving an obstruction theory for the horizontal map by \cite[3.15]{Manolache2012Pullbacks}. 
% We have $$\frak M_{0,2}^{ss} = \cl T \times B\bb G_m\, ,$$
% and the map $B \bb G_m \to \on{Spec} \field$ is unobstructed. Giving 
% an obstruction theory for the diagonal map is again equivalent to
% giving an obstruction theory for the horizontal map of 
% \begin{equation}\label{hhww33}
%  \begin{tikzcd}
%   \cat{Rub}_{g,A}(X, \ca L)' \arrow[dr] \arrow[r] & \frak M_{g,n} \times \cl T \times B\bb G_m \arrow[d]\\
% &\frak M_{g,n} \times \cl T\, . 
% \end{tikzcd}
% \end{equation}
% If we start with the Marcus-Wise obstruction theory for the diagonal map
% of \eqref{hhww33}, we obtain an obstruction theory for the horizontal map of
% \eqref{hhww33}. The latter is the diagonal map of \eqref{kk882}, so we obtain
% an obstruction theory for the horizontal map of \eqref{kk882}.
% The virtual fundamental classes of
% $\cat{Rub}_{g,A}(X, \ca L)'$
% associated to the Marcus-Wise theory and the induced obstruction theory
% of the horizontal map of \eqref{kk882} are equal by \cite{Manolache2012Pullbacks}. 


% Graber and Vakil \cite{Graber2005Relative-virtua} give a pointwise description of a perfect relative obstruction theory for the map
% \begin{equation}
% \cat{Rub}_{g,A}(X, \ca L)' \to \MbarX' \times \frak M_{0,2}^{ss}. 
% \end{equation}
% To make the comparison to \cite{Marcus2017Logarithmic-com}, we first construct a relative obstruction theory for this map in the setting of \cite{Marcus2017Logarithmic-com}. Then we easily see that it gives the same virtual class as the Gysin pullback, and then verify by a `pointwise' computation that this matches up with the steps given in \cite{Graber2005Relative-virtua}. %\{Not sure what exactly to say here... GV only describe the graded pieces of the obstruction theory pointwise, so we can't actually prove that what we write down is the same. What we do is the same as the approach Marcus and Wise took (assuming I understand it correctly). }

% We factor the map 
% \begin{equation}
% \cat{Rub}_{g,A}(X, \ca L)' \to \MbarX' \times \frak M_{0,2}^{ss} 
% \end{equation}
% as
% \begin{equation}\label{eq:key_factorisation}
% \cat{Rub}_{g,A}(X, \ca L)' \to \MbarX' \times_{\frak M_{g,n}} \cat{Rub}_{g,A} \to \MbarX' \times \frak M_{0,2}^{ss}\,. 
% \end{equation}
% The first map is the top left vertical arrow in \eqref{eq:long_vertical_pullback}, hence has a perfect relative obstruction theory given by $R^1\pi_*\ca O_C$. To obtain an obstruction theory for the second map we consider the pullback diagram 
% \begin{equation}
%  \begin{tikzcd}
% \MbarX' \times_{\frak M_{g,n}} \cat{Rub}_{g,A} \arrow[r] \arrow[d] & \cat{Rub}_{g,A}  \arrow[d]  \\
% \MbarX'\times \cl T \arrow[r] & \frak M_{g,n} \times \cl T . 
% \end{tikzcd}
% \end{equation}
% %\begin{equation}
% % \begin{tikzcd}
% %\MbarX' \times_{\frak M_{g,n}} \cat{Rub}_{g,A} \arrow[r] \arrow[d] & \cat{Rub}_{g,A} \times B\bb G_m \arrow[d] & \\
% %\MbarX'\times \frak M_{0,2}^{ss} \arrow[r] & \frak M_{g,n} \times \frak M_{0,2}^{ss} \arrow[r, equals] &  \frak M_{g,n} \times \cl T \times B\bb G_m. 
% %\end{tikzcd}
% %\end{equation}
% %\todo{Before I wrote: `here something not quite right; really want $\MbarX'\times \ca T$ at bottom left, to make cartesian. Not really a problem as $B\bb G_m$ is smooth, but need to get correct. ' But now with absolute version things seem to work better? }
% Now, for the right-hand vertical map an obstruction is described by \cite[Proposition \MWref{5.6.5.3}\MWvone{5.20}]{Marcus2017Logarithmic-com}, so we get an obstruction for the left vertical map by base-change. \{AAAAAAh! It's $\cl T$ vs $\frak M_{0,2}^{ss}$ again! }

% Summarising these obstruction theories, we have one step given by $R^1\pi_* \ca O_C$, and the other given terms of data at nodes, coming from \cite[Proposition \MWref{5.6.5.3}\MWvone{5.20}]{Marcus2017Logarithmic-com}. As this obstruction theory arises by base-change from one constructed just above, it still gives the same virtual class as the Gysin pullback. It remains to compare to \cite{Graber2005Relative-virtua}. 


\subsubsection{Marcus-Wise and Graber-Vakil}\label{mwgv}
\source{pixtons-formula-abel-jacobi-picard-stack}


As recalled above, Marcus-Wise define a two-step obstruction theory for the map 
\begin{equation*}
\cat{Rub}_{g,A}(X, \ca L)' \to \frak M_{g,n} \times \cl T\,. 
\end{equation*}
% with the factorisation 
% $$\cat{Rub}_{g,A}(X, \ca L)' \to \cat{Rub}_{g,A} \to  \frak M_{g,n} \times \cl T \, .$$
Graber and Vakil consider an obstruction theory for the map 
\begin{equation} \label{vrrt}
\source{pixtons-formula-abel-jacobi-picard-stack}
\cat{Rub}_{g,A}(X, \ca L)' \to \MbarX' \times \frak M_{0,2}^{ss}\, .
\end{equation}

We wish to show an equality of the corresponding virtual fundamental classes on $\cat{Rub}_{g,A}(X, \ca L)'$. Since $$\frak M_{0,2}^{ss} = \cl T \times B\bb G_m\, ,$$ and that the maps $\MbarX' \to \frak M_{g,n}$ and $B \bb G_m \to \on{Spec} \field$ are unobstructed, we have an unobstructed map
$$ \MbarX' \times \frak M_{0,2}^{ss} \to \frak M_{g,n} \times \frak M_{0,2}^{ss} \to \frak M_{g,n} \times \cl T\, .$$
Our final step is therefore to compare the obstruction theories (and thereby the corresponding virtual pullbacks) between Marcus-Wise and Graber-Vakil \cite{Graber2005Relative-virtua,Li2002A-degeneration-}.
%The final step is to show that the obstruction theory induced by Marcus-Wise
%for the map \eqref{vrrt} matches the obstruction theory studied by Graber-Vakil
%in \cite{Graber2005Relative-virtua}.
%The latter is obtained from \cite{Li2002A-degeneration-}.
We will match the obstruction spaces when the base $S$ is a point.
The full matching of deformation theories is similar and will be treated in \cite{PW}.
The claims are also required for \cite{Marcus2017Logarithmic-com}.

Suppose we are given the data of a point in $\cat{Rub}_{g,A}(X, \ca L)'(S)$,
\begin{equation}
(\tau\colon \tilde C \to C_S, \R/X_S, \phi\colon \tilde C \to \R, f\colon C_S \to X_S, p_1, \dots, p_n, f^*\ca L \isom \ca O_C(\alpha))\, .
\end{equation}

\begin{proposition}
\source{pixtons-formula-abel-jacobi-picard-stack}
\notready\label{prop:vfc_comparison_summary}
\source{pixtons-formula-abel-jacobi-picard-stack}
The restriction to $\MbarX'$ of the class $\DR_{g,A}(X, \ca L)$ of \cite{Janda2018Double-ramifica} is equal to the class obtained by letting $\DRop_{g,A}$ act on the fundamental class of $\MbarX'$ via the map induced by $\ca L$. 
\end{proposition}



%The full deformation theory will
%However, there is an obstruction to proving this coming from the fact that Graber and Vakil only describe the graded pieces of their obstruction theories, and only in the case where the base $S$ is a point (or more generally of constant degeneration), see \cref{rem:pointwise}. The same obstruction is encountered in \cite{Marcus2017Logarithmic-com}. 

%To facilitate the comparison to \cite{Marcus2017Logarithmic-com} w
%We recall some details of the two-step obstruction theory of Graber and Vakil. As this is a pointwise theory the description only makes sense when the base $S$ is a geometric point (or more generally when the curves are of constant degeneration). 
%\todo{Y. Is it really true that [MW] compared two obstruction theories over a geometric point?}\{I think so. I think that's what the word `pointwise' means on first line of Section 5.6, page 18 of [MW]. Note also that prop 5.12 of [MW] seems only to make sense over a point, or at least for curves of constant degeneration. I think the reason for this is that for formulae given in [GV] for the obstruction groups (with the $f^\dagger$ etc also only make sense over a point, so in general it would not be clear what [MW] should be trying to compare to. }

\begin{proof}
The primary obstruction of Graber-Vakil lies in 
\begin{equation}\label{ggtt2}
\source{pixtons-formula-abel-jacobi-picard-stack}
%H^0(\tilde C, r^{-1}\ca{E}xt^1(\Omega_{R/X_S}(\log D), \ca O_{R}))\, .
H^0(\tilde C, \phi^{-1}\ca{E}xt^1(\Omega_{\R/X_S}(\log D), \ca O_{\R}))\, .
\end{equation}
Here, $D$ is the divisor on $\R$ given by the sum of the two markings $r_0$ and $r_\infty$, and $\Omega_{\R/X_S}(\log D)$ is the sheaf of relative 1-forms on $\R/X_S$ allowed logarithmic poles along $D$ (a coherent sheaf on $\R$). The obstruction space 
\eqref{ggtt2}
is isomorphic to the product of the deformation spaces of the nodes of $\R$
and coincides with the obstruction space
for the map $$\widetilde{\cat{Rub}}_{g,A} \to \frak M_{g,n} \times \cl T$$ the vertical arrow in \eqref{eq:obs_factor_MW}), coming from \cite[Proposition \MWref{5.6.5.3}\MWvone{5.20}]{Marcus2017Logarithmic-com} (where they assume the divisibility conditions of Lemma \ref{lem:divisibility}, hence the results apply to $\widetilde{\cat{Rub}}_{g, A}$ and not to $\cat{Rub}_{g, A}$). 

Suppose that the primary obstruction vanishes. Denote by $$T_{\R/X_S} = \sheafhom_{\ca O_{\R}}(\Omega_{\R/X_S} , \ca O_{\R})$$ the relative tangent sheaf. There is a secondary obstruction in 
\begin{equation}\label{pp39}
\source{pixtons-formula-abel-jacobi-picard-stack}
H^1(\tilde{C}, \phi^{\dagger}(T_{\R/X_S}))\,, 
\end{equation}
where the $\phi^{\dagger}$ is the torsion-free part of $\phi^*$, see \cite{Graber2005Relative-virtua} and \cite{Marcus2017Logarithmic-com}. The
obstruction space \eqref{pp39}
is the image of the obstruction space $H^1 (C_S, f^* T_\R ) = H^1(C_S, \ca O_{C_S})$ for the map $$\cat{Rub}_{g,A}(X, \ca L)' \to \widetilde{\cat{Rub}}_{g,A}\, $$
the horizontal arrow in \eqref{eq:obs_factor_MW}, coming from \cite[Proposition \MWref{5.6.3.1}\MWvone{5.17}]{Marcus2017Logarithmic-com}. 



%\jocomment{R: The way the logic goes, we have to connect the GV obstructions to
%he induced theory of 6.4.3. You are matching (88). Could you put some
%explicit sentences (which are now implicit) with the connections to 6.4.3?}
%\{The references you point to have changed a bit, but I think the question still stands; to what extent does comparing these cohomology groups somehow `match up' the obstruction theories? Interpreting `obstruction spaces' to be the graded pieces of the two-step theories, I guess it matches these up as you wrote a couple of paras above. But I do not think we can really say more than that at this point. The only data GV give us is the two graded pieces of their theory in the case when the base $S$ is a point. Following [MW] I write that GV give a two-step theory, and that is consistent with them having a `primary and secondary obstruction', but I don't even know what the implicit factorisation is (maybe Jonathan does), let alone how to describe the actual obstruction theory. What we do is the same as [MW]...
%}

%\begin{remark}\label{rem:pointwise}\{Perhaps we don't want to include this paragraph, but I found it helpful to understand. }
%Note that formation of the relative tangent sheaf $$T_{\R/X_S} = \sheafhom_{\ca O_{\R}}(\Omega_{\R/X_S} , \ca O_{\R})$$ does not commute with non-flat base-change over $S$ unless $[R/\bb G_m]$ is constant over $S$. The simplest case where this can be seen is when $X$ is a point, $S$ is a trait with uniformiser $t$, and $R$ is locally given by $\ca O_S[x,y]/(xy-t)$. Then $\sheafhom_{\ca O_{\R}}(\Omega_{\R/X_S} , \ca O_{\R})$ is a reflexive sheaf on $R$, and $R$ is regular hence reflexive implies locally free. On the other hand, over the closed point $(t)$ the sheaf $\sheafhom_{\ca O_{\R}}(\Omega_{\R/X_S} , \ca O_{\R})$ is torsion-free but not locally free. This justifies our assertion that Graber and Vakil's description of the graded pieces of the obstruction theory is only valid pointwise. 
%\end{remark}



When both of these obstructions vanish, the deformations are a torsor under $H^0(\tilde{C}, \phi^\dagger T_{\R/X_S} )$, an extension of the first term of the obstruction complex in \cite[Proposition \MWref{5.6.5.3}\MWvone{5.20}]{Marcus2017Logarithmic-com} by the first term of the obstruction complex of \cite[Proposition \MWref{5.6.4.1}\MWvone{5.17}]{Marcus2017Logarithmic-com}. The comparison of the obstruction theories is complete. %We have proven the following result.
\end{proof}

%
%\section{OLDER STUFF, to go}
%
%%[I'm not sure how hard the general case will be. Marcus and Wise do it when the target is a point, and it's pretty painful, but maybe for general targets one can piggyback on their work. Hopefully we anyway don't need to do this. ]
%
%%Remember the goal is to prove \cref{lem:projective_space_comparison}. 
%
%
%%For this case this should be a lot easier; we have $X = \bb P^k$, and $\ca L = \ca O(1)$, and we're interested in the locus where the pullback to $C$ is relatively (very) ample, the space of stable maps is smooth, and there will be less virtualisation. 
%
%%Hope: on both side the obstruction will just vanish. Then the VFC is just the fundamental class, and the stacks are the same, so we are done. I (David) am currently checking this on the Log/Div side. Maybe someone could check it on the rubber maps side? 
%
%
%%Idea: in 5.5 of Marcus and Wise, can I just take $X = \Mbar(X, \beta)$, $C$ the universal marked curve, and $L$ the pullback of the bundle on $X$? But their comparison is only for the case of the trivial bundle (that's what GV did, it appears)? 
%
%%Can we just use [MW, Theorem B] (obstruction is $H^1(C, \ca O_C)$?! [Later: No, that's the obstruction OVER Rub, not absolutie/over Mbar].
%
%% We have a map 
%%\begin{equation}
%%\Mbar(X, \beta) \to \frak{Pic}; (C, p_i, f) \mapsto f^*\ca O(1)(-AP). 
%%\end{equation}
%%But I'm confused; in the Div/Rub setting we never virtualise anyway! Well, implicitly we do; we take a gysin pullback. And this is (as we proved carefully in the Jussieu paper) the same as taking the VFC where the obstruction is pulled back from the inclusion of Div into Pic (actually we proved something slightly less general, but same argument works here). 
%%
%%So what we want to prove is \\
%%1.  the stack of rubber maps coincides with the naive pullback of Rub as a stack; \\
%%2. The obstruction theory on rubber maps coincides with the pull back of the obstruction on Rub. 
%
%%Now, the pulled-back obstruction on Rub is probably just $H^1(C, \ca O_C)$, if I'm reading [MW THeorem B] correctly? Maybe prthe same is true on the rubber maps side?!
%
%%===
%
%I think we just have to follow through Section 5 of [MW] with the following replacements: 
%\begin{enumerate}
%\item
%$\ca M_{g,n}^{log}$ replaced by $\Mbar_{g,n}(X, \beta)$ (with pullback log structure from $\frak M$). 
%\item $\ca O_C$ replaced by $f^* \ca L$ on the universal curve over $\Mbar_{g,n}(X, \beta)$. 
%\end{enumerate}
%%Now the obstructions described by [GV05] in this case seem basically the same; they don't seem sensitive to these changes. I hope/expect that on the MW side things will go through the same. At a first look it's not clear where these changes will actually make a difference. Maybe the answer is `nowhere, really'? In some places MW work with $\omega^k$, which we can avoid, but this seems not relevant for the comparison to [GV05]. 
%
%%p.s. not $100\%$ sure on this. Can't we subsume the $f^* \ca L$ into the $\R$ bundle (in MW terminology)? Need to understand better the role of this $\R$. Later: No, we can't. $\R$ lives on $S$, which for us would be $\Mbar(X, \beta)$. Whereas our line bundle $\ca L$ lives on $X$ itself. So we will need some modification to the theory in MW; but hopefully it won't be too bad. 
%
%
%
%%Let $X = \bb P^k$, and $\ca L = \ca O_X(1)$. Then $\MbarX$ is a smooth DM stack.\textcolor{brown}{Y. Maybe we can remove this sentence?} \{Yes, I think that's fine. And probably we should, else it's not actually a smooth DM stack, right? Though I'm not sure just now where we use smoothness below... } We also have fixed markings $A$. Then w
%We have maps
%\begin{equation*}
%\MbarX \to \frak Pic; \;\;\; (C, P, f) \mapsto f^*\ca L. 
%\end{equation*}
%and 
%\begin{equation*}
%\MbarX \to \frak Pic; \;\;\; (C, P, f) \mapsto f^*\ca L (-AP). 
%\end{equation*}
%
%We can form pullback diagrams
%\begin{equation}
% \begin{tikzcd}
%\frak M_A(\ca L) = \cat{Rub}_{g, A}(\MbarX, f^*\ca L)  \arrow[r] \arrow[d] & \MbarX \arrow[d, "f^*\ca L"] \\
% \cat{Rub}_{g,A} \arrow[r] & \frak Pic
%\end{tikzcd}
%\end{equation}
%and
%\begin{equation}
% \begin{tikzcd}
%\frak M(\ca L(-AP)) =  \cat{Rub}_{g, 0}(\MbarX, f^*\ca L(-AP))  \arrow[r] \arrow[d] & \MbarX \arrow[d, "f^*\ca L(-AP)"] \\
% \cat{Rub}_{g,0} \arrow[r] & \frak Pic
%\end{tikzcd}
%\end{equation}
%which are just the same up to an automorphism of $\frak Pic$ given by tensoring with $\ca O(AP)$. Note $\cat{Rub}_{g,0}$  is just the pullback of $\cat{Rub}_g$ along the smooth map that forgets the markings. 
%
%Then the universal DR cycle is by definition a virtual fundamental class of $ \cat{Rub}_{g, 0}(\MbarX, f^*\ca L(-AP))$, but to translate into the language of [MW, Section 5.6] we want to look instead at $ \cat{Rub}_{g, A}(\MbarX, f^*\ca L)$, which comes to the same thing as they just differ by a translation on $\frak Pic$. 
%
%Now, we defined $\DRop$ as an operational Chow class. We have a relative perfect obstruction theory for $s : \cat{Rub}_{g,A} \to \frak Pic$ and by the pullback property of relative perfect obstruction theory \cite[Proposition 7.2]{Behrend1997The-intrinsic-n} we obtain a relative perfect obstruction theory on $ \cat{Rub}_{g, A}(\MbarX, f^*\ca L)$, which corresponds to the operational Chow class $\DRop$ defined by the unit section, cf. \cite{D_Jussieu_paper_near_end}.  
%
%Now consider the map 
%\begin{equation}
%\frak M_A(\ca L) = \cat{Rub}_{g, A}(\MbarX, f^*\ca L) \to \MbarX \times_\Lambda \frak M^{ss}_{0,2}
%\end{equation}
%sending $(\tau\colon \tilde C \to C_S, Q/S, \phi\colon \tilde C \to Q, f\colon C \to X_S, p_1, \dots, p_n)$ to the pair consisting of 
%\begin{enumerate}
%\item
%the stabilisation of $C\to X_S$; 
%\item the curve $Q$. 
%\end{enumerate}
%\textcolor{brown}{(Y. The following argument on relative perfect obstruction theory can be justified on the locus $\MbarX'$. Consider the morphism
%\[
%\frak{M}_{A}(\ca L)' \to \MbarX' \times_\Lambda \frak M^{ss}_{0,2} \to \frak M_{g,n} \times_\Lambda \frak M^{ss}_{0,2}.
%\]
%Since $\MbarX'$ is unobstructed, we get a canonical relative perfect obstruction theory for the first morphism by \cite[Remark 3.15]{Manolache2012Pullbacks}. I think this is the thing described below.)} \{Can you explain what are you taking as $\frak M$, $F$ and $G$ in Manolache's notation. I assume $F =\frak{M}_{A}(\ca L)'$, and I guess we must take $G =  \MbarX' \times_\Lambda \frak M^{ss}_{0,2}$, but I'm not sure what $\frak M$ wants to be. }
%
%Then [GV05] give a pointwise perfect relative obstruction theory for this map. To facilitate the comparison we recall some details of this two-step obstruction theory. As this is a pointwise theory the description only makes sense when the base $S$ is a geometric point (or more generally when the curves are of constant degeneration). 
%
%Suppose we are given 
%\begin{equation}
%(\tau\colon \tilde C \to C_S, Q/S, \phi\colon \tilde C \to Q, f\colon C \to X_S, p_1, \dots, p_n) \in  \cat{Rub}_{g, A}(\MbarX, f^*\ca L)
%\end{equation}
%We define a semistable curve $R/X_S$ by $R = \ca L^\vee \otimes Q \in \frak M_{0,2}^{ss}(X_S)$, which comes with a map $r\colon \tilde C \to R$ induced by the maps $\tilde C \to Q$ and the (unfixed?!) isomorphism $f^*\ca L \isom \ca O_C(\alpha)$ (need to add this into the above data). 
%
%
%The primary obstruction lies in 
%\begin{equation}
%H^0(\tilde C, r^{-1}\ca{E}xt^1(\Omega^{log}_{R/X_S}, \ca O_{R})), 
%\end{equation}
%which is isomorphic to the product of the deformation spaces of the nodes of $R$, see [MW Proposition 5.6.5.3]; the latter is proven under a stability assumption but this assumption is not used in the proof. 
%%does not really use stability as far as I can see. Note that the formula above is the same as 
%%\begin{equation}
%%H^0(\tilde C, ?^{-1}\ca{E}xt^1(\Omega^{log}_{Q/S}, \ca O_{Q})), 
%%\end{equation}
%%where $?$ is the as-yet-unnamed map $\tilde C \to Q$, since they differ just by pulling back along a smooth map then tensoring with a line bundle. I'm not sure which is really `naturally the obstruction'; they're both just the spaces of local deformations of nodes. 
%
%Suppose that the primary obstruction vanished. We consider the log canonical bundle $\Omega^{log}_{R/X_S}$, a coherent sheaf on $R$, and the log relative tangent bundle $T^{log}_{R/X_S} = \sheafhom_{\ca O_{R}}(\Omega^{log}_{R/X_S} , \ca O_{R})$, then there is a secondary obstruction in 
%\begin{equation}
%H^1(C, r^{\dagger}(T^{log}_{R/X_S})), 
%\end{equation}
%where the $r^{\dagger}$ is the torsion-free part of $r^*$, cf. GV or MW. This is reminiscent of Proposition 5.6.4.1 of [MW], but is not the same because of the dagger (see later in [loc.cit.]. 
%
%%But this is the `secondary obstruction'; first there is an obstruction coming from the nodes, which needs to vanish before this comes into play. 
%
%%Note: Also, this is all very weird, talking about `nodes' and `torsion part'. This only makes sense when $S$ is a point, or more generally $Q$ is constant, otherwise these notions will not commute with base-change. I guess this is what Marcus and Wise talk about with a `point wise obstruction theory'. This is troubling. 
%
%%The primary obstruction: I think this is 
%%\begin{equation}
%%H^0(\tilde C, r^{-1}\ca{E}xt^1(\Omega^{log}_{R/X_S}, \ca O_{R})), 
%%\end{equation}
%%which is isomorphic to the product of the deformation spaces of the nodes of $R$, see [MW Proposition 5.6.5.3]; the latter does not really use stability as far as I can see. Note that the formula above is the same as 
%%\begin{equation}
%%H^0(\tilde C, ?^{-1}\ca{E}xt^1(\Omega^{log}_{Q/S}, \ca O_{Q})), 
%%\end{equation}
%%where $?$ is the as-yet-unnamed map $\tilde C \to Q$, since they differ just by pulling back along a smooth map then tensoring with a line bundle. I'm not sure which is really `naturally the obstruction'; they're both just the spaces of local deformations of nodes. 
%
%Hopefully we have correctly described GV's `point wise obstruction theory' in the language of MW. Still need to check that it actually coincides with the obstruction theory of MW. I think this is not actually a problem. In 5.5 and 5.6 we've just made appropriate changes to work relative to $X$, nothing much else changes. The heart of the comparison is the computations in 5.6.5 and 5.6.6. Now in 5.6.5 the target is never relevant, and the same arguments seem (to me) to go through fine in the prestable case (though they are only stated in the stable case). Then I think the discussion above roughly translates 5.6.6 into the `target X' setting. OK, I should sleep on this and see if I still feel optimistic in the morning... 
%
%
%
%\subsection{Comparing Rub and Div}
%This just uses that they are birational. I think Rosa is looking into this. 
%
